\documentclass[10pt]{amsart}
\usepackage[a4paper,margin=22mm]{geometry}
\usepackage[no-math]{fontspec}
\setmainfont{Latin Modern Roman}[SmallCapsFont={lmromancaps10-regular.otf}]
\usepackage{amsmath,amssymb,amsthm,mathtools,hyperref,xcolor,bm,array}
\usepackage[all,cmtip]{xy}
\usepackage[shortlabels]{enumitem}
\usepackage[numbers,square]{natbib}
\newtheorem{theorem}{Theorem}[section]
\newtheorem{lemma}[theorem]{Lemma}
\newtheorem{proposition}[theorem]{Proposition}
\newtheorem{corollary}[theorem]{Corollary}
\theoremstyle{definition}
\newtheorem{defn}[theorem]{Definition}
\newcommand{\ie}{i.e.\ }
\newcommand{\eg}{e.g.\ }
\newcommand{\resp}{respectively\ }
\emergencystretch=3em
\newcommand{\comaj}{{\mathrm {comaj}}}
\newcommand{\maj}{{\mathrm {maj}}}
\newcommand{\inv}{{\mathrm {inv}}}
\newcommand{\minimaj}{{\mathrm {minimaj}}}
\newcommand{\sign}{{\mathrm {sign}}}
\newcommand{\area}{{\mathrm {area}}}
\newcommand{\dinv}{{\mathrm {dinv}}}
\newcommand{\cont}{{\mathrm {cont}}}
\newcommand{\mult}{{\mathrm {mult}}}
\newcommand{\Des}{{\mathrm {Des}}}
\newcommand{\Val}{{\mathrm {Val}}}
\newcommand{\Rise}{{\mathrm {Rise}}}
\newcommand{\Stir}{{\mathrm {Stir}}}
\newcommand{\Hilb}{{\mathrm {Hilb}}}
\newcommand{\grFrob}{{\mathrm {grFrob}}}
\newcommand{\des}{{\mathrm {des}}}
\newcommand{\asc}{{\mathrm {asc}}}
\newcommand{\coinv}{{\mathrm {coinv}}}
\newcommand{\rev}{{\mathrm {rev}}}
\newcommand{\Asc}{{\mathrm {Asc}}}
\newcommand{\Coinv}{{\mathrm {Coinv}}}
%
\newcommand{\Ind}{{\mathrm {Ind}}}
\newcommand{\Res}{{\mathrm {Res}}}
%
\newcommand{\NSym}{{\mathbf {NSym}}}
\newcommand{\Sym}{{\mathrm {Sym}}}
\newcommand{\QSym}{{\mathrm {QSym}}}
\newcommand{\SYT}{{\mathrm {SYT}}}
\newcommand{\SSK}{{\mathrm {SSK}}}
\newcommand{\Frob}{{\mathrm {Frob}}}
\newcommand{\triv}{{\mathrm {triv}}}
\newcommand{\init}{{\mathrm {in}}}
\newcommand{\Ch}{{\mathrm {Ch}}}
\newcommand{\ch}{{\mathbf {ch}}}
%
\newcommand{\neglex}{{\mathtt {neglex}}}
\newcommand{\iDes}{{\mathrm {iDes}}}
\newcommand{\shape}{{\mathrm {shape}}}
%
\newcommand{\symm}{{\mathfrak{S}}}
\newcommand{\wt}{{\mathrm{wt}}}
%
\newcommand{\CC}{{\mathbb {C}}}
\newcommand{\QQ}{{\mathbb {Q}}}
\newcommand{\ZZ}{{\mathbb {Z}}}
\newcommand{\NN}{{\mathbb{N}}}
\newcommand{\RR}{{\mathbb{R}}}
\newcommand{\OP}{{\mathcal{OP}}}
\newcommand{\FF}{{\mathbb{F}}}
\newcommand{\RRR}{{\mathcal{R}}}
\newcommand{\KK}{{\mathbb{K}}}
\newcommand{\WWW}{{\mathcal{W}}}
\newcommand{\NNN}{{\mathcal{N}}}
\newcommand{\FFF}{{\mathcal{F}}}
\newcommand{\CCC}{{\mathcal{C}}}
\newcommand{\LD}{{\mathcal {LD}}}
\newcommand{\AAA}{{\mathcal{A}}}
\newcommand{\MMM}{{\mathcal{M}}}
\newcommand{\BBB}{{\mathcal{B}}}
\newcommand{\GS}{{\mathcal{GS}}}
%
\newcommand{\zz}{{\mathbf {z}}}
\newcommand{\xx}{{\mathbf {x}}}
\newcommand{\II}{{\mathbf {I}}}
\newcommand{\yy}{{\mathbf {y}}}
\newcommand{\TT}{{\mathbf {T}}}
\newcommand{\mm}{{\mathbf {m}}}
\newcommand{\ii}{{\mathbf{i}}}
\newcommand{\sss}{{\mathbf{s}}}
\newcommand{\dd}{{\mathbf {d}}}
\newcommand{\LM}{{\textsc {lm}}}
%
\newcommand{\pib}{\overline{\pi}}
\newcommand{\edit}[1]{{\color{red}{\bf {#1}}}}

\title[Projective zero-Hecke decompositions]{Projective zero-Hecke decompositions of modified ordered-set-partition coinvariant quotients}
\author{Jia Huang; Brendon Rhoades}
\date{}
\begin{document}
\maketitle
\noindent\textbf{Source.} Original Theorem 5.9, Ordered set partitions and the 0-Hecke algebra, Algebraic Combinatorics 1 (2018), 47--80, DOI 10.5802/alco.10. Original algebra/quotient/composition definitions, point-set and Gr\"obner dimension argument, descent-monomial straightening and Demazure basis proofs, and the complete module decomposition proof follow. Unrelated symmetric-function characteristic theory and illustrative Young diagrams are outside this unit; their custom style is not loaded. All selected mathematical macros are explicitly preserved from the native preamble or supplied by genuine AMS/xy/bm runtime packages.
\section{Introduction}
\label{Introduction}
The symmetric group $\symm_n$ acts on the polynomial ring
$\QQ[\xx_n] := \QQ[x_1, \dots, x_n]$ by variable permutation.
The corresponding \emph{invariant subring} $\QQ[\xx_n]^{\symm_n}$ consists of all $f \in \QQ[\xx_n]$ with $w(f) = f$  
for all $w \in \symm_n$, 
and is generated 
by the elementary symmetric functions $e_1(\xx_n), \dots, e_n(\xx_n)$, where
\begin{equation}
e_d(\xx_n) = e_d(x_1, \dots, x_n) := \sum_{1 \leq i_1 < \cdots < i_d \leq n} x_{i_1} \cdots x_{i_d}.
\end{equation}
The \emph{invariant ideal} $I_n \subseteq \QQ[\xx_n]$ is the ideal generated by those invariants
$\QQ[\xx_n]^{\symm_n}_+$ with
vanishing constant term:
\begin{equation}
I_n := \langle \QQ[\xx_n]^{\symm_n}_+ \rangle = \langle e_1(\xx_n), \dots, e_n(\xx_n) \rangle.
\end{equation}
The \emph{coinvariant algebra} $R_n := {\QQ[\xx_n]} / {I_n}$ is the corresponding quotient ring. 

Let $k \leq n$ be two positive integers.  Haglund, Rhoades, and Shimozono \cite[Defn. 1.1]{HRS} introduced
the ideal $I_{n,k} \subseteq \QQ[\xx_n]$ with generators 
\begin{equation}
I_{n,k} := \langle x_1^k, x_2^k, \dots, x_n^k, e_n(\xx_n), e_{n-1}(\xx_n), \dots, e_{n-k+1}(\xx_n) \rangle
\end{equation}
and studied the corresponding quotient ring $R_{n,k} := {\QQ[\xx_n]} / {I_{n,k}}$.  Since $I_{n,k}$ is homogeneous
and
stable 
under the action of $\symm_n$, the ring $R_{n,k}$ is a graded $\symm_n$-module.
When $k = n$, we have $I_{n,n} = I_n$, so that   $R_{n,n} = R_n$ and we recover the usual 
invariant ideal and coinvariant algebra.

To study $R_{n,k}$ one needs the notion of an \emph{ordered set partition} of $[n] := \{1, 2, \dots, n\}$, which is a set partition of $[n]$ with a total order on its blocks.
For example, we have an ordered set partition
\begin{equation*}
\sigma = (25 \mid 6 \mid 134)
\end{equation*}
written in the `bar notation'. The three blocks $\{2,5\}$, $\{6\}$, and $\{1,3,4\}$ are ordered from left to right, and elements of each block are increasing.

Let $\OP_{n,k}$ denote the collection of ordered set partitions of $[n]$ with $k$ blocks.  We have
\begin{equation}
|\OP_{n,k}| = k! \cdot \Stir(n,k),
\end{equation}
where $\Stir(n,k)$ is the (signless) Stirling number of the second kind counting $k$-block
set partitions of $[n]$.
The symmetric group $\symm_n$ acts on $\OP_{n,k}$ by permuting the letters $1,\ldots,n$.  For example, the permutation $w=241365$, written in one-line notation, sends $(25 \mid 6 \mid 134)$ to $(46 \mid 5 \mid 123)$.
Now let $\FF$ be an arbitrary field and let $n$ be a positive integer.  
The 
\emph{(type A) $0$-Hecke algebra} $H_n(0)$ 
is the unital associative $\FF$-algebra with generators $\pi_1, \pi_2, \dots, \pi_{n-1}$ and relations
\begin{equation}
\label{zero-hecke-relations}
\begin{cases}
\pi_i^2 =  \pi_i & 1 \leq i \leq n-1, \\
\pi_i \pi_j = \pi_j \pi_i & |i - j| > 1,  \\
\pi_i \pi_{i+1} \pi_i = \pi_{i+1} \pi_i \pi_{i+1} & 1 \leq i \leq n-2.
\end{cases}
\end{equation}
%We will find it convenient to consider an alternative set of generators $\pib_1, \pib_2, \dots, \pib_{n-1}$
%for $H_n(0)$, where $\pib_i := \pi_i - 1$.  The generators $\pib_i$ satisfy the relations
%\begin{equation}
%\begin{cases}
%\pib_i^2 =  -\pib_i & 1 \leq i \leq n-1, \\
%\pib_i \pib_j = \pib_j \pib_i & |i - j| > 1,  \\
%\pib_i \pib_{i+1} \pib_i = \pib_{i+1} \pib_i \pib_{i+1} & 1 \leq i \leq n-2.
%\end{cases}
%\end{equation}
Recall that the symmetric group $\symm_n$ has Coxeter generators $\{s_1, s_2, \dots, s_{n-1} \}$, where
$s_i$ is the adjacent transposition $s_i = (i, i+1)$.
These generators satisfy similar relations as \eqref{zero-hecke-relations} except that $s_i^2=1$ for all $i$. 
If $w \in \symm_n$ is a permutation and $w = s_{i_1} \cdots s_{i_{\ell}}$ is a reduced (i.e., as short as possible)
expression for $w$
in the Coxeter generators $\{s_1, \dots, s_{n-1}\}$, we define 
the 0-Hecke algebra element $\pi_w := \pi_{i_1} \cdots \pi_{i_{\ell}} \in H_n(0)$.  It can be shown that the 
set $\{ \pi_w \,:\, w \in \symm_n \}$ forms a basis for $H_n(0)$ as an $\FF$-vector space, and in particular
$H_n(0)$ has dimension $n!$. 
In contrast to the situation with the symmetric group, the representation theory of the 
0-Hecke algebra is insensitive to the choice of ground field, which motivates our generalization
from $\QQ$ to $\FF$.


The algebra $H_n(0)$ is a deformation of the symmetric group algebra
$\FF[\symm_n]$.  Roughly speaking, whereas in a typical $\FF[\symm_n]$-module the 
generator $s_i$ acts by `swapping' the letters $i$ and $i+1$, in a typical 
$H_n(0)$-module the generator $\pi_i$ acts by `sorting' the letters $i$ and $i+1$.
Indeed, the relations satisfied by the $\pi_i$ are precisely the relations satisfied by bubble sorting operators
acting on a length $n$ list of entries $x_1 \dots x_n$  from a totally ordered alphabet:
\begin{equation}
\pi_i .( x_1 \dots x_i x_{i+1} \dots x_n) := \begin{cases}
x_1 \dots x_{i+1} x_i \dots x_n & x_i > x_{i+1} \\
x_1 \dots x_i x_{i+1} \dots x_n & x_i \leq x_{i+1}.
\end{cases}
\end{equation}


Proving 0-Hecke analogs of  module theoretic results concerning the symmetric group has received
a great deal of recent study in algebraic combinatorics \cite{BBSSZ, Huang, HuangTab, VT}; let us 
recall the 0-Hecke analog of the variable permutation action of $\symm_n$ on a polynomial ring.

Let $\FF[\xx_n] := \FF[x_1, \dots, x_n]$ be the polynomial ring in $n$ variables over the field $\FF$.
The algebra $H_n(0)$ acts on $\FF[\xx_n]$ by the \emph{isobaric Demazure operators}:
\begin{equation}\label{Demazure}
\pi_i(f) := \frac{x_i f - x_{i+1} (s_i(f))}{x_i - x_{i+1}},\quad 1\le i\le n-1.
\end{equation}
% where $s_i.f$ is the polynomial obtained by interchanging $x_i \leftrightarrow x_{i+1}$ in $f$.
% (We already mentioned the permutation action w.f in the beginning of this section.)
If $f \in \FF[\xx_n]$ is symmetric in the variables $x_i$ and $x_{i+1}$, then $s_i(f) = f$ and thus $\pi_i(f) = f$.
The isobaric Demazure operators give a 0-Hecke analog of variable permutation.


We also have a 0-Hecke analog of the permutation action of $\symm_n$ on $\OP_{n,k}$.
It is well-known that the 0-Hecke algebra $H_n(0)$ has another generating set $\{\pib_1,\ldots,\pib_{n-1}\}$ subject to the relations
\begin{equation}\label{pi-bar-relations}
\begin{cases}
\pib_i^2 =  -\pib_i & 1 \leq i \leq n-1, \\
\pib_i \pib_j = \pib_j \pib_i & |i - j| > 1,  \\
\pib_i \pib_{i+1} \pib_i = \pib_{i+1} \pib_i \pib_{i+1} & 1 \leq i \leq n-2.
\end{cases}
\end{equation}
Here $\pib_i:=\pi_i-1$ for all $i$.
We will often use the relation $\pib_i\pi_i = \pi_i\pib_i=0$.
One can define $\pib_w:=\pib_{i_1}\cdots\pib_{i_\ell}$ for any $w\in\symm_n$ with a reduced expression $w=s_{i_1}\cdots s_{i_\ell}$ and show that the set $\{\pib_w:w\in\symm_n\}$ is a basis for $H_n(0)$.
Let $\FF[\OP_{n,k}]$ be the $\FF$-vector space with basis given by $\OP_{n,k}$.
Then $H_n(0)$ acts on $\FF[\OP_{n,k}]$ by the rule
\begin{equation}
\label{pi-defining-action}
\pib_i . \sigma := \begin{cases}
-\sigma, & \text{if $i+1$ appears in a block to the left of $i$ in $\sigma$,} \\
s_i(\sigma), & \text{if $i+1$ appears in a block to the right of $i$ in $\sigma$,} \\
0, & \text{if $i+1$ appears in the same block as $i$ in $\sigma$,} 
%\pi_i . \sigma := \begin{cases}
%0 & \text{if $i+1$ appears in a block to the left of $i$ in $\sigma$,} \\
%\sigma + s_i(\sigma), & \text{if $i+1$ appears in a block to the right of $i$ in $\sigma$.} \\
%\sigma & \text{if $i+1$ appears in the same block as $i$ in $\sigma$,} 
\end{cases}
\end{equation}
% where $s_i.\sigma$ is the ordered set partition $\sigma$ with $i$ and $i+1$ interchanged.
% This is already defined.
For example, we have
\begin{align*}
\pib_1 (25 \mid 6 \mid 134) &= -(25 \mid 6 \mid 134),  \\
\pib_2 (25 \mid 6 \mid 134) &= (35 \mid 6 \mid 124), \\
\pib_3 (25 \mid 6 \mid 134) &= 0.
%\pi_1 (25 \mid 6 \mid 134) &= 0, \\
%\pi_2 (25 \mid 6 \mid 134) &= (25 \mid 6 \mid 134) + (35 \mid 6 \mid 124), \\
%\pi_3 (25 \mid 6 \mid 134) &= (25 \mid 6 \mid 134).
\end{align*}
It is straightforward to check that these operators satisfy the relations \eqref{pi-bar-relations}
and so define an $H_n(0)$-action on $\FF[\OP_{n,k}]$. 
In fact, this is a special case of an $H_n(0)$-action on generalized ribbon tableaux introduced in~\cite{HuangTab}.
See also the proof of Lemma~\ref{osp-substructure}.



The coinvariant algebra $R_n$ can be viewed as a 0-Hecke module.
Indeed, 
%we may express the operator $\pi_i$ as a composition $\pi_i = \partial_i x_i$, where $\partial_i$ is the \emph{divided difference operator}
%\begin{equation}
%\partial_i(f) := \frac{f - s_i(f)}{x_i - x_{i+1}}.
%\end{equation}
the ``Leibniz rule''
%Brendon: We have $\pib_i(f g) = \pib_i(f) g + s_i(f) \pib_i(g)$ for $\pib_i:=\pi_i-1$, exactly the same as the Leibniz rule for $\partial_i$.
%To avoid using $\pib_i$ in this section, we need a slighly modified rule: $\pi_i(fg) = \pi_i(f) g + s_i(f) (\pi_i(g)-g)$. 
%Either way seems fine to me. --Jia})
\begin{equation}
\label{leibniz}
\pib_i(fg) = \pib_i(f) g + s_i(f) \pib_i(g)
%\pi_i(fg) = \pi_i(f) g + s_i(f) (\pi_i(g)-g)
\end{equation}
implies that the ideal 
$I_n \subseteq \FF[\xx_n]$ generated by $e_1(\xx_n), \dots, e_n(\xx_n) \in \FF[\xx_n]$ 
is stable under the action of $H_n(0)$ on $\FF[\xx_n]$.  Therefore,
 the quotient $R_n = {\FF[\xx_n]} / {I_n}$ inherits a 0-Hecke action.
We will study a 0-Hecke analog of the rings $R_{n,k}$ of Haglund, Rhoades, and Shimozono \cite{HRS}.
For $k < n$ the ideal $I_{n,k}$ is not usually stable under the action of $H_n(0)$ on 
$\FF[\xx_n]$, so that the quotient ring $R_{n,k} = {\FF[\xx_n]} / {I_{n,k}}$ does not have the structure of 
an $H_n(0)$-module.  To remedy this situation, we introduce the following modified family of ideals.
Let 
\begin{equation}
h_d(x_1, \dots, x_i) := \sum_{1 \leq j_1 \leq \cdots \leq j_d \leq i} x_{j_1} \cdots x_{j_d} 
\end{equation}
be the complete homogeneous symmetric function of degree $d$ in the variables $x_1, x_2, \dots, x_i$.

\begin{defn}
For two positive integers $k \leq n$, we define a quotient ring 
\begin{equation*}
S_{n,k} := {\FF[\xx_n]} / {J_{n,k}}
\end{equation*}
where $J_{n,k} \subseteq \FF[\xx_n]$ is the ideal with generators
\begin{equation*}
J_{n,k} := \langle h_k(x_1), h_k(x_1, x_2), \dots, h_k(x_1, x_2, \dots, x_n), e_n(\xx_n), e_{n-1}(\xx_n), \dots, e_{n-k+1}(\xx_n) 
\rangle.
\end{equation*}
\end{defn}


The ideal $J_{n,k}$ is homogeneous.
We claim that $J_{n,k}$ is stable under the action of $H_n(0)$.  Since $e_d(\xx_n) \in \FF[\xx_n]^{\symm_n}$
and $h_k(x_1, \dots, x_i)$ is symmetric in $x_j$ and $x_{j+1}$ for $j \neq i$, 
thanks to Equation~\eqref{leibniz} this reduces to the observation that
\begin{equation}\label{shift}
\pi_i(h_k(x_1, \dots, x_i)) = h_k(x_1, \dots, x_i, x_{i+1}).
\end{equation}
Equation \eqref{shift} is clear when $i=1$ and can be obtained from the following identity when $i\ge2$:
\begin{equation}
h_k(x_1,\ldots,x_i) = \sum_{0\le j\le k} x_i^j h_{k-j}(x_1,\ldots,x_{i-1}). 
\end{equation}
%More details below: 
%Write $h_k^{(i)} := h_k(x_1,\ldots,x_i)$. We have
%\[ \pi_1(h_k^{(1)}) = \pi_1(x_1^k) = x_1^k + x_1^{k-1}x_2 + \cdots + x_1x_2^{k-1}+x_2^k = h_k^{(2)}.\]
%For $i\ge2$ we have
%\[ h_k^{(i)} = \sum_{0\le j\le k} x_i^j h_{k-j}^{(i-1)}. \]
%Thus 
%\begin{eqnarray*}
%\pi_i(h_k^{(i)}) & = & 
%\frac{x_i h_k^{(i)} - x_{i+1} s_i(h_k^{(i)}) } { x_i - x_{i+1} } \\
%&=& \sum_{0\le j\le k} \frac{ (x_i^{j+1}  - x_{i+1}^{j+1}) h_{k-j}^{(i-1)} } { x_i - x_{i+1} } \\
%&=& \sum_{0\le j\le k} \sum_{0\le \ell\le j} x_i^\ell x_{i+1}^{j-\ell} h_{k-j}^{(i-1)} \\
%&=& h_k^{(i+1)}.
%\end{eqnarray*}
Thus the quotient
$S_{n,k}$ has the structure of a graded $H_n(0)$-module.


It can be shown that $J_{n,n} = I_n$, so that $S_{n,n} = R_n$ is the classical coinvariant module.  
At the other extreme, we have $J_{n,1} = \langle x_1, x_2, \dots, x_n \rangle$, so that 
$S_{n,1} \cong \FF$ is the trivial $H_n(0)$-module in degree $0$.
\section{Background}
\label{Background}


\subsection{Compositions}
Let $n$ be a nonnegative integer.
A \emph{(strong) composition} $\alpha$ of $n$ is a sequence
$\alpha = (\alpha_1, \dots, \alpha_{\ell})$ of positive integers with $\alpha_1 + \cdots + \alpha_{\ell} = n$.
We call $\alpha_1,\ldots,\alpha_\ell$ the \emph{parts} of $\alpha$.
We write $\alpha \models n$ to mean that $\alpha$ is a  composition of $n$.
We also write $|\alpha| = n$ for the \emph{size} of $\alpha$ and $\ell(\alpha) = \ell$ for the number  of parts of $\alpha$.


The \emph{descent set} $\Des(\alpha)$ of a composition $\alpha = (\alpha_1, \dots, \alpha_{\ell}) \models n$ is the subset of $[n-1]$ given by
\begin{equation}
\Des(\alpha) := \{\alpha_1, \alpha_1 + \alpha_2, \dots, \alpha_1 + \alpha_2 + \cdots + \alpha_{\ell - 1} \}.
\end{equation}
The map $\alpha \mapsto \Des(\alpha)$ gives a bijection from the set of compositions of $n$ to the collection  
of subsets of $[n-1]$.
The \emph{major index} of $\alpha = (\alpha_1, \dots, \alpha_{\ell})$
is 
\begin{equation}
\maj(\alpha) := \sum_{i\in\Des(\alpha)} i = (\ell - 1) \cdot \alpha_1 + \cdots + 1 \cdot \alpha_{\ell-1} + 0 \cdot \alpha_{\ell}.
\end{equation}
Given two compositions $\alpha, \beta \models n$, we write $\alpha \preceq \beta$ if 
$\Des(\alpha) \subseteq \Des(\beta)$.  Equivalently, we have $\alpha \preceq \beta$ if the composition $\alpha$
can be formed by merging adjacent parts of the composition $\beta$.  
If $\alpha \models n$, the \emph{complement} $\alpha^c \models n$ of $\alpha$ is the unique composition of $n$
which satisfies $\Des(\alpha^c) = [n-1] \setminus \Des(\alpha)$.


As an example of these concepts, let $\alpha = (2,3,1,2) \models 8$.  We have $\ell(\alpha) = 4$. 
The descent set of $\alpha$ is $\Des(\alpha) = \{2,5,6\}$.
The major index is $\maj(\alpha) = 2+5+6 = 3 \cdot 2 + 2 \cdot 3 + 1 \cdot 1 + 0 \cdot 2 = 13$.
The complement of $\alpha$ is $\alpha^c = (1,2,1,3,1) \models 8$ with descent set $\Des(\alpha^c) = \{1,3,4,7\} = [7]\setminus\{2,5,6\}$.  



If $\ii = (i_1, \dots, i_n)$ is any sequence of  integers, the \emph{descent set}
$\Des(\ii)$ is given by
\begin{equation}
\Des(\ii) := \{1 \leq j \leq n-1 \,:\, i_j > i_{j+1} \}.
\end{equation}
The \emph{descent number} of $\ii$ is $\des(\ii) := |\Des(\ii)|$ and the \emph{major index} of $\ii$ is $\maj(\ii) := \sum_{j \in \Des(\ii)} j$.
Finally, the \emph{inversion number} $\inv(\ii)$ 
\begin{equation}
\inv(\ii) := | \{ (j,j') : 1\le j<j'\le n,\ i_j > i_{j'} \}|
\end{equation}
counts the number of inversion pairs in the sequence $\ii$.

If a permutation $w \in \symm_n$ has one-line notation $w = w(1) \cdots w(n)$, we define 
$\Des(w)$, $\maj(w)$, $\des(w)$, and $\inv(w)$ as in the previous paragraph for the sequence $(w(1),\ldots,w(n))$.
It turns out that $\inv(w)$ is equal to the \emph{Coxeter length} $\ell(w)$ of $w$, i.e., the length of
a reduced expression for $w$ in the generating set $\{s_1, \dots, s_{n-1}\}$ of 
$\symm_n$.
Moreover, we have $i\in\Des(w)$ if and only if some reduced expression of $w$ ends with $s_i$.
We also let $\iDes(w) := \Des(w^{-1})$ be the descent set of the inverse of the permutation $w$.

The statistics $\maj$ and $\inv$ are equidistributed on $\symm_n$ and their common distribution has
a nice form.    Let us recall the standard $q$-analogs of numbers, factorials, and multinomial coefficients:
\begin{align*}
[n]_q := 1 + q + \cdots + q^{n-1} &   &[n]!_q := [n]_q [n-1]_q \cdots [1]_q  \\
{n \brack a_1, \dots , a_r}_q := \frac{[n]!_q}{[a_1]!_q \cdots [a_r]!_q} 
& &{n \brack a}_q := \frac{[n]!_q}{[a]!_q [n-a]!_q}.
\end{align*}
MacMahon \cite{MacMahon} proved  
\begin{equation}
\sum_{w \in \symm_n} q^{\inv(w)} = \sum_{w \in \symm_n} q^{\maj(w)} = [n]!_q,
\end{equation}
and any statistic on $\symm_n$ which shares this distribution is called \emph{Mahonian}.
The joint distribution $\sum_{w \in \symm_n} q^{\inv(w)} t^{\maj(w)}$ of the pair of statistics
$(\inv, \maj)$ is called the \emph{biMahonian distribution}.




If $\alpha \models n$ and $\ii = (i_1, \dots, i_n)$ 
is a sequence of  integers of length $n$, we define $\alpha \cup \ii \models n$ 
to be the unique composition of $n$ which satisfies
\begin{equation}
\Des(\alpha \cup \ii) = \Des(\alpha) \cup \Des(\ii).
\end{equation}
For example, let $\alpha = (3,2,3) \models 8$ and let
$\ii = (4,5,0,0,1,0,2,2)$.  We have 
\begin{equation*}
\Des(\alpha \cup \ii) = \Des(\alpha) \cup \Des(\ii) = \{3,5\} \cup \{2,5\} = \{2,3,5\},
\end{equation*}
so that $\alpha \cup \ii = (2,1,2,3)$.
Whenever $\alpha \models n$ and $\ii$ is a length $n$ sequence, we have the 
relation $\alpha \preceq \alpha \cup \ii$.

A \emph{partition} $\lambda$ of $n$ is a weakly decreasing sequence $\lambda = (\lambda_1 \geq \cdots \geq \lambda_{\ell})$
of positive integers
which satisfies $\lambda_1 + \cdots + \lambda_{\ell} = n$.  We write
$\lambda \vdash n$ to mean that $\lambda$ is a partition of $n$.  We also write
$|\lambda| = n$ for the \emph{size} of $\lambda$ and $\ell(\lambda) = \ell$ for the number of \emph{parts}
of $\lambda$.
The \emph{(English) Ferrers diagram}
of $\lambda$ consists of $\lambda_i$ left justified boxes in row $i$.




Identifying partitions with Ferrers diagrams, if $\mu \subseteq \lambda$ are a pair of partitions related by 
containment, the \emph{skew partition} $\lambda/\mu$ is obtained by removing $\mu$ from $\lambda$.
We write $|\lambda/\mu| := |\lambda| - |\mu|$ for the number of boxes in this skew diagram.
A \emph{ribbon} is an edgewise connected skew diagram which contains no $2 \times 2$ square.
The set of compositions of $n$ is in bijective correspondence with the set of size $n$ ribbons: a composition
$\alpha = (\alpha_1, \dots, \alpha_{\ell})$ corresponds to the ribbon whose $i^{th}$ row from the bottom
contains $\alpha_i$ boxes.  We will identify compositions with ribbons in this way.
Let $\alpha \models n$ be a composition.  We define a permutation $w_0(\alpha) \in \symm_n$ as follows.
Starting at the leftmost column and working towards the right, and moving from top to bottom within each column, fill the 
ribbon diagram of $\alpha$ with the numbers $1, 2, \dots, n$ (giving a standard tableau).
The permutation $w_0(\alpha)$ has one-line notation obtained by reading along the ribbon from the bottom row to the top row, proceeding from left to right within each row.  
It can be shown that $w_0(\alpha)$ is the unique left weak Bruhat minimal permutation $w \in \symm_n$ which
satisfies $\Des(w) = \Des(\alpha)$ (cf. Bj\"orner and Wachs~\cite{BjornerWachs}).
\subsection{Ordered set partitions}
As explained in Section~\ref{Introduction},
an \emph{ordered set partition} $\sigma$ of size $n$ is a set partition of $[n]$ with a total order on its blocks.
Let $\OP_{n,k}$ denote the collection of ordered set 
partitions of size $n$ with $k$ blocks.  
In particular, we may identify $\OP_{n,n}$ with $\symm_n$.

Also as in Section~\ref{Introduction},
we write an ordered set partition of $[n]$ as a permutation of $[n]$ with bars to separate blocks, such that letters within each block are increasing and blocks are ordered from left to right.
For example, we have
\begin{equation*}
\sigma = (245 \mid 6 \mid 13) \in \OP_{6,3}.
\end{equation*}
The \emph{shape} of an ordered set partition $\sigma = (B_1 \mid \cdots \mid B_k)$ is
the composition $\alpha = (|B_1|, \dots, |B_k|)$.  For example, the above ordered set partition has shape $(3,1,2) \models 6$.

If $\alpha \models n$ is a composition, let $\OP_{\alpha}$ denote the collection of ordered set partitions of $n$ with shape $\alpha$.  
Given an ordered set partition $\sigma \in \OP_{\alpha}$, we can also represent $\sigma$ as the pair
$(w, \alpha)$, where $w = w(1) \cdots w(n)$ is the permutation in $\symm_n$ (in one-line notation) obtained by erasing 
the bars in $\sigma$.  For example, the above ordered set partition becomes
\begin{equation*}
\sigma = (245613, (3,1,2)).
\end{equation*}
This notation establishes a bijection between $\OP_{n,k}$ and pairs $(w, \alpha)$ where $\alpha \models n$ 
is a composition with $\ell(\alpha) = k$ and $w \in \symm_n$ is a permutation with $\Des(w) \subseteq \Des(\alpha)$.


%There are two different generalizations of Mahonian statistics from $\symm_n$ to $\OP_{n,k}$ which will come up
%in our work.  The first of these distributions is given by the following extension of 
%$\inv$.  Let $\sigma \in \OP_{n,k}$ be an ordered set partition represented as a pair $(w, \alpha)$ as above.
%We define
%\begin{equation}
%\inv(\sigma) = \inv(w,\alpha) := \inv(w).
%\end{equation}
%For example, we have
%\begin{equation*}
%\inv(24 \mid 3 \mid 15) = \inv(24315) = 4.
%\end{equation*}
%If $\alpha = (\alpha_1, \dots, \alpha_{\ell}) \models n$,
%the statistic $\inv$ has distribution on $\OP_{\alpha}$ given by the $q$-multinomial coefficient:
%\begin{equation}
%\sum_{\sigma \in \OP_{\alpha}} q^{\inv(\sigma)} = {n \brack \alpha_1, \dots, \alpha_{\ell}}_q.
%\end{equation}
%This is the usual distribution of Coxeter length on minimal parabolic coset representatives in $\symm_n$.
%By summing this formula over all compositions $\alpha \models n$ with $\ell(\alpha) = k$, we get the 
%distribution of $\inv$ on $\OP_{n,k}$.



We extend the statistic $\maj$ from permutations to ordered set partitions as follows.
Let $\sigma = (B_1 \mid \cdots \mid B_k) \in \OP_{n,k}$ be an ordered set partition represented as a pair $(w, \alpha)$ as above.  
We define the \emph{major index} $\maj(\sigma)$ to be the statistic
\begin{equation}
\maj(\sigma) = \maj(w, \alpha) := \maj(w) + \sum_{i \,  : \,  \max(B_i) < \min(B_{i+1})} (\alpha_1 + \cdots + \alpha_i - i).
\end{equation}
For example, if $\sigma = ( 2 4 \mid 5 7 \mid 1 3 6 \mid 8)$, then
\begin{equation*}
\maj(\sigma) = \maj(24571368) + (2 - 1) + (2 + 2 + 3 - 3) = 4 + 1 + 4 = 9.
\end{equation*}

We caution the reader that our definition of
$\maj$ is \emph{not} equivalent to, or even equidistributed with,
the corresponding statistics for ordered set partitions in \cite{RW, HRS} and elsewhere.  
However,
the distribution of our $\maj$ on $\OP_{n,k}$ is the reversal of the distribution of their $\maj$.


The generating function for $\maj$ on $\OP_{n,k}$ may be described as follows.  
Let $\rev_q$ be the operator on polynomials in the variable $q$ which reverses coefficient sequences.
For example, we have
\begin{equation*}
\rev_q(3q^3 + 2q^2 + 1) = q^3 + 2q + 3.
\end{equation*}
The \emph{$q$-Stirling number} $\Stir_q(n,k)$ is defined by the recursion
\begin{equation}
\Stir_q(n,k) = \Stir_q(n-1,k-1) + [k]_q \cdot \Stir_q(n-1,k)
\end{equation}
and the initial condition $\Stir_q(0,k) = \delta_{0,k}$, where $\delta$ is the Kronecker delta.

\begin{proposition}
Let $k \leq n$ be positive integers.  We have
\begin{equation}
\sum_{\sigma \in \OP_{n,k}} q^{\maj(\sigma)} = \rev_q( [k]!_q \cdot \Stir_q(n,k)).
\end{equation}
\end{proposition}

\begin{proof}
To see why this equation holds, consider the statistic $\maj'$ on an ordered set partition 
$\sigma = (B_1 \mid \cdots \mid B_k) = (w, \alpha) \in \OP_{n,k}$ defined by
\begin{equation}
\maj'(\sigma) = \maj'(w,\alpha) := \sum_{i = 1}^k (i-1)(\alpha_i - 1) + \sum_{i \, : \, \min(B_i) > \max(B_{i+1})} i.
\end{equation}
This is precisely the version of major index on ordered set partitions studied
by Remmel and Wilson \cite{RW}.  They proved \cite[Eqn. 15, Prop. 5.1.1]{RW}
that
\begin{equation}
\label{remmel-wilson}
\sum_{\sigma \in \OP_{n,k}} q^{\maj'(\sigma)} = [k]!_q \cdot \Stir_q(n,k).
\end{equation}

On the other hand, for any $\sigma = (w,\alpha) = (B_1 \mid \cdots \mid B_k) \in \OP_{n,k}$ we have
\begin{equation}
\maj(w) = \sum_{i \, : \, \max(B_i) > \min(B_{i+1})} (\alpha_1 + \cdots + \alpha_i).
\end{equation}
This implies
\begin{align}
\maj(\sigma) &= \maj(w) + \sum_{i \, : \, \max(B_i) < \min(B_{i+1})} (\alpha_1 + \cdots + \alpha_i - i) \\
&= \sum_{i = 1}^{k-1} [(k-i) \cdot \alpha_i]   - \sum_{i \, : \, \max(B_i) < \min(B_{i+1})} i.
\end{align}

The longest element $w_0 =  n \dots 21$ (in one-line notation) of $\symm_n$ gives an involution on $\OP_\alpha$ by
\begin{equation*}
\sigma = (B_1 \mid \cdots \mid B_k) \mapsto w_0(\sigma) = (w_0(B_1) \mid \cdots \mid w_0(B_k)).
\end{equation*}



If $\alpha \models n$ and $\ell(\alpha) = k$, then for $\sigma = (B_1 \mid \cdots \mid B_k) \in \OP_{\alpha}$
and any index $1 \leq i \leq k-1$ 
we have $\max(B_i) < \min(B_{i+1})$ if and only if $\min(w_0(B_i)) > \max(w_0(B_{i+1}))$.
Therefore,
\begin{align}
\maj'(\sigma) + \maj(w_0(\sigma)) &=  \sum_{i = 1}^{k} [(i-1)(\alpha_i - 1) + (k-i) \cdot \alpha_i] \\
&=  \sum_{i = 1}^{k} [-\alpha_i - i + 1 + k \alpha_i] \\
&= (k-1)(n - k) + {k \choose 2}.
\end{align}
On the other hand, it is easy to see that 
\begin{equation*}
\max\{ \maj(\sigma) \,:\, \sigma \in \OP_{n,k} \} = (k-1)(n - k) - {k \choose 2} =
\max\{ \maj'(\sigma) \,:\, \sigma \in \OP_{n,k} \}.
\end{equation*}
Applying Equation~\eqref{remmel-wilson} gives
\begin{equation}
\sum_{\sigma \in \OP_{n,k}} q^{\maj(\sigma)} = \rev_q \left[ \sum_{\sigma \in \OP_{n,k}} q^{\maj'(\sigma)} \right]
= \rev_q( [k]!_q \cdot \Stir_q(n,k) ).
\end{equation}
\end{proof}




We have an action of the $0$-Hecke algebra $H_n(0)$ on $\FF[\OP_{n,k}]$ given
by Equation~\eqref{pi-defining-action}.
%For $i\in[n-1]$ and $\sigma\in\OP_{n,k}$ let
%\begin{equation}
%\pib_i(\sigma) := 
%\begin{cases}
%-\sigma & \text{if $i+1$ appears in a block to the left of $i$ in $\sigma$,} \\
%0 & \text{if $i+1$ appears in the same block as $i$ is $\sigma$, and} \\
%s_i(\sigma) & \text{if $i+1$ appears in a block to the right of $i$ in $\sigma$.}
%\end{cases}
%\end{equation}
This $H_n(0)$-action preserves $\FF[\OP_\alpha]$ for each composition $\alpha$ of $n$.
%see the proof of Proposition~\ref{osp-substructure}.




\subsection{Gr\"obner theory}
We review material related to Gr\"obner bases of ideals $I \subseteq \FF[\xx_n]$ and standard monomial
bases of the corresponding quotients $\FF[\xx_n]/I$.  For a more leisurely introduction to this material,
see \cite{CLO}.


A total order $\leq$ on the monomials in $\FF[\xx_n]$ is called a \emph{term order} if 
\begin{itemize}
\item $1 \leq m$ for all monomials $m \in \FF[\xx_n]$, and
\item  if $m, m', m'' \in \FF[\xx_n]$ are monomials, then $m \leq m'$ implies $m \cdot m'' \leq m' \cdot m''$.
\end{itemize}
In this paper, we will consider the lexicographic term order \emph{with respect to the variable ordering
$x_n > \cdots > x_2 > x_1$}.  That is, we have
\begin{equation*}
x_1^{a_1} \cdots x_n^{a_n} < x_1^{b_1}  \cdots x_n^{b_n}
\end{equation*}
if and only if there exists an integer $1 \leq  j \leq n$ such that $a_{j+1} = b_{j+1}, \dots, a_n = b_n$, and
$a_j < b_j$.
Following the notation of $\textsc{sage}$,
we call this term order {\tt neglex}.

Let $\leq$ be any term order on monomials in $\FF[\xx_n]$.  If $f \in \FF[\xx_n]$ is a nonzero polynomial,
let $\init_<(f)$ be the leading (i.e., largest) term of $f$ with respect to $<$.  If $I \subseteq \FF[\xx_n]$ is an ideal, the 
associated \emph{initial ideal} is the monomial ideal
\begin{equation}
\init_<(I) := \langle \init_<(f) \,:\, f \in I - \{0\} \rangle.
\end{equation}
The set of monomials
\begin{equation}
\{ \text{monomials $m \in \FF[\xx_n]$} \,:\, m \notin \init_<(I) \}
\end{equation}
descends to a $\FF$-basis for the quotient  ${\FF[\xx_n]} / {I}$; this basis is called the 
\emph{standard monomial basis} (with respect to the term order $\leq$) \cite[Prop. 1, p. 230]{CLO}.

Let $I \subseteq \FF[\xx_n]$ be any ideal and let $\leq$ be a term order. 
A finite set $G = \{g_1, \dots, g_r\} \subseteq I$ 
of nonzero polynomials in $I$ is called a \emph{Gr\"obner basis} of $I$ if
\begin{equation}
\init_<(I) = \langle \init_<(g_1), \dots , \init_<(g_r) \rangle.
\end{equation}
If $G$ is a Gr\"obner basis of $I$, then we have $I = \langle G \rangle$ \cite[Cor. 6, p. 77]{CLO}.



Let $G$ be a Gr\"obner basis for $I$ with respect to the
term order $\leq$.  The basis $G$ is called \emph{minimal} if
\begin{itemize}
\item for any $g \in G$, the leading coefficient of $g$ with respect to $\leq$ is $1$, and
\item for any $g \neq g'$ in $G$, the leading monomial of $g$ does not divide the leading monomial of $g'$.
\end{itemize}
A minimal Gr\"obner basis $G$ is called \emph{reduced} if in addition 
\begin{itemize}
\item  for any $g \neq g'$ in $G$, the leading monomial of $g$ does not divide any term in the polynomial $g'$.
\end{itemize}
Up to a choice of term order,
every ideal $I$ has a unique reduced Gr\"obner basis \cite[Prop. 6, p. 92]{CLO}.



Now let us recall the 0-Hecke analog of the above story.
Consider an arbitrary ground field $\FF$.
The representation theory of the $\FF$-algebra $H_n(0)$ was studied by Norton \cite{Norton}
and has a different flavor from that of $\QQ[\symm_n]$ since $H_n(0)$ is not semisimple.
Norton \cite{Norton} proved that %the set of algebra elements
%\begin{equation}
%\{ \pib_{w_0(\alpha)} \pi_{w_0(\alpha^c)} \,:\, \alpha \models n \}
%\end{equation}
% is, up to scaling, a complete set of primitive orthogonal idempotents for $H_n(0)$. 
% Unfortunately, the elements $\pib_{w_0(\alpha)} \pi_{w_0(\alpha^c)}$ are neither orthogonal nor idempotent.
% A set of primitive orthongal idempotents for $H_n(0)$ is given by 
% Tom Denton, {\it A Combinatorial Formula for Orthogonal Idempotents in the 0-Hecke Algebra of the Symmetric Group}.
the $H_n(0)$-modules 
\begin{equation}
P_{\alpha} := H_n(0) \pib_{w_0(\alpha)} \pi_{w_0(\alpha^c)},
\end{equation}
where $\alpha$ ranges over all compositions of $n$, form a complete list of nonisomorphic
 indecomposable projective 
$H_n(0)$-modules.  
For each $\alpha\models n$, the $H_n(0)$-module $P_\alpha$ has a basis 
\[ \{ \pib_w \pi_{w_0(\alpha^c)} : w\in\symm_n,\ \Des(w) = \Des(\alpha) \}. \]
Moreover, $P_\alpha$ has a unique maximal submodule spanned by all elements in the above basis except its cyclic generator $\pib_{w_0(\alpha)} \pi_{w_0(\alpha^c)}$, and the quotient of $P_\alpha$ by this maximal submodule, denoted by $C_\alpha$, is one-dimensional and admits an $H_n(0)$-action by $\pib_i = -1$ for all $i\in \Des(\alpha)$ and $\pib_i = 0$ for all $i\in \Des(\alpha^c)$. 
The collections $\{ P_\alpha:\alpha\models n\}$ and $\{ C_\alpha: \alpha\models n\}$ are complete lists of nonisomorphic projective indecomposable and irreducible $H_n(0)$-modules, respectively.

\section{Hilbert Series and Artin basis}
\label{Hilbert}


\subsection{The point sets $Z_{n,k}$}
In this section we will prove that $\dim(S_{n,k}) = |\OP_{n,k}|$.
To do this, we will use tools from elementary algebraic geometry.
This basic method dates back to the work of Garsia and Procesi on Springer
fibers and Tanisaki quotients \cite{GP}.

Given a finite point set $Z \subseteq \FF^n$, let $\II(Z) \subseteq \FF[\xx_n]$ be the ideal of 
polynomials which vanish on $Z$:
\begin{equation}
\II(Z) := \{ f \in \FF[\xx_n] \,:\, f(\zz) = 0 \text{ for all $\zz \in Z$} \}.
\end{equation} 
There is a natural identification of the quotient ${\FF[\xx_n]} / {\II(Z)}$ with the collection of
polynomial functions $Z \rightarrow \FF$.


We claim that any function $Z \rightarrow \FF$ may be realized as the restriction of a polynomial
function.  This essentially follows from Lagrange Interpolation.  Indeed, since $Z \subseteq \FF^n$ is finite, there exist
field elements $\alpha_1, \dots, \alpha_m \in \FF$ such that $Z \subseteq \{\alpha_1, \dots, \alpha_m\}^n$. 
For any $n$-tuple of integers $(i_1, \dots, i_n)$ between $1$ and $m$, the polynomial
\begin{equation*}
\prod_{j_1 \neq i_1} (x_1 - \alpha_{j_1}) \cdots \prod_{j_n \neq i_n} (x_n - \alpha_{j_n}) \in \FF[\xx_n]
\end{equation*} 
vanishes on every point of $\{\alpha_1, \dots, \alpha_m\}^n$ except for 
$(\alpha_{i_1}, \dots, \alpha_{i_n})$.  Hence, an arbitrary $\FF$-valued function on
$\{\alpha_1, \dots, \alpha_m\}^n$ may be realized using a linear combination of polynomials of the above form.
Since $Z \subseteq \{\alpha_1, \dots, \alpha_m\}^n$, the same is true for an arbitrary $\FF$-valued function on $Z$.

By the last paragraph, we may identify the quotient ${\FF[\xx_n]} / {\II(Z)}$ with the collection of all 
functions $Z \rightarrow \FF$.
In particular, the dimension of this quotient as an $\FF$-vector space is
\begin{equation}
\dim \left( {\FF[\xx_n]} / {\II(Z)} \right) = |Z|.
\end{equation}

The ideal $\II(Z)$ is almost never homogeneous.  To get a homogeneous ideal, we do the following.
For any nonzero polynomial $f \in \FF[\xx_n]$, let $\tau(f)$ be the top degree component of $f$.  That 
is, if $f = f_d + f_{d-1} + \cdots + f_0$ where $f_i$ has homogeneous degree $i$ for all $i$ and $f_d \neq 0$, then
$\tau(f) = f_d$.  The ideal $\TT(Z) \subseteq \FF[\xx_n]$ is generated by the top degree components of all
nonzero polynomials in $\II(Z)$.  In symbols:
\begin{equation}
\TT(Z) := \langle \tau(f) \,:\, f \in \II(Z) - \{0\} \rangle.
\end{equation}

The ideal $\TT(Z)$ is homogeneous by definition, so that ${\FF[\xx_n]} / {\TT(Z)}$ is a graded $\FF$-vector space.
Moreover, it is well known that
\begin{equation}
\dim \left( {\FF[\xx_n]} / {\TT(Z)} \right) = \dim \left( {\FF[\xx_n]} / {\II(Z)} \right) = |Z|.
\end{equation}


Our three-step strategy for proving $\dim(S_{n,k}) = |\OP_{n,k}|$ is as follows.

\begin{enumerate}
\item  Find a  point set $Z_{n,k} \subseteq \FF^n$ which is in bijective correspondence with $\OP_{n,k}$.
\item  Prove that the generators of $J_{n,k}$ arise as top degree components of polynomials in $\II(Z_{n,k})$,
so that $J_{n,k} \subseteq \TT(Z_{n,k})$.
\item Use Gr\"obner theory to prove $\dim(S_{n,k}) \leq |\OP_{n,k}|$, forcing $\dim(S_{n,k}) = |\OP_{n,k}|$
by Steps 1 and 2.
\end{enumerate}

A similar three-step strategy was used by Haglund, Rhoades, and Shimozono \cite{HRS} 
in their analysis of the  $\symm_n$-module structure of $R_{n,k}$.
In our setting, since we do not have a group action, we can only use this strategy to deduce the  vector
space structure of $S_{n,k}$, rather than the $H_n(0)$-module structure of $S_{n,k}$.

To achieve Step 1 of our strategy, we need to find a candidate set $Z_{n,k} \subseteq \FF^n$ which is in bijective
correspondence with $\OP_{n,k}$.  Here we run into a problem:  to define our candidate point sets, we need the 
field $\FF$ to contain at least $n+k-1$ elements.  This problem did not arise in the work of Haglund et. al. \cite{HRS};
they worked exclusively over the field $\QQ$.  To get around this problem, we use the following  trick.

\begin{lemma}
\label{algebraic-swindle}
Let $\FF \subseteq \KK$ be a field extension and $J = \langle f_1, \dots, f_r \rangle \subseteq \FF[\xx_n]$ an ideal of $\FF[\xx_n]$ generated by $f_1, \dots, f_r \in \FF[\xx_n]$.  
Then $\dim_\FF ( {\FF[\xx_n]} / {J} )  = \dim_\KK ( {\KK[\xx_n]} / J' )$ where $J' := {\KK\otimes_\FF J}$.
\end{lemma}

Since $J_{n,k}$ is generated by polynomials with all coefficients equal to $1$, the generating set of $J_{n,k}$ 
satisfies the conditions of Lemma~\ref{algebraic-swindle}.

\begin{proof}
Let $\leq$ be any term order.
It suffices to show that the quotient rings ${\FF[\xx_n]} / {J}$ and ${\KK[\xx_n]} / {J'}$ have the same 
standard monomial bases with respect to $\leq$.  
To calculate the reduced Gr\"obner basis for the ideal $J$, we apply Buchberger's Algorithm~\cite[Ch. 2, \S7]{CLO} 
to the generating set $\{f_1, \dots, f_r\}$.   
To calculate the Gr\"obner basis for the ideal $J'$, we also apply Buchberger's Algorithm to the generating set
$\{f_1, \dots, f_r\}$.  In either case, all of the coefficients involved in the polynomial long division are contained in the
field $\FF$.  In particular, the reduced Gr\"obner bases of $J$ and $J'$ coincide.  Hence, the standard monomial
bases of ${\FF[\xx_n]} / {J}$ and ${\KK[\xx_n]} / {J'}$ also coincide.
\end{proof}

We are ready to define our point sets $Z_{n,k}$.  Thanks to Lemma~\ref{algebraic-swindle}, we may harmlessly
assume that 
the field $\FF$ has at least $n+k-1$ elements by replacing $\FF$ with an extension if necessary.
We will have to choose a somewhat non-obvious point set $Z_{n,k} \subseteq \FF^n$ in order to get the 
desired equality of ideals $\TT(Z_{n,k}) = J_{n,k}$.

\begin{defn}
Assume $\FF$ has at least $n+k-1$ elements and let $\alpha_1, \alpha_2, \dots, \alpha_{n+k-1} \in \FF$
be a list of $n+k-1$ distinct field elements.  Define $Z_{n,k} \subseteq \FF^n$ to be the collection of points
$(z_1, z_2, \dots, z_n)$ such that
\begin{itemize}
\item  for $1 \leq i \leq n$ we have $z_i \in \{\alpha_1, \alpha_2, \dots, \alpha_{k+i-1} \}$,
\item the coordinates $z_1, z_2, \dots, z_n$ are distinct, and
\item  we have $\{\alpha_1, \alpha_2, \dots, \alpha_k\} \subseteq \{z_1, z_2, \dots, z_n\}$.
\end{itemize}
\end{defn}

We claim that $Z_{n,k}$ is in bijective correspondence with $\OP_{n,k}$.  A bijection
$\varphi: \OP_{n,k} \rightarrow Z_{n,k}$ may be obtained as follows.
Let $\sigma = (B_1 \mid \cdots \mid B_k) \in \OP_{n,k}$ be an ordered set partition;
we define $\varphi(\sigma) = (z_1, \dots, z_n) \in Z_{n,k}$.
For $1 \leq i \leq k$, we first set $z_j = \alpha_i$, where $j = \min(B_i)$.  
Write the set of unassigned indices of
$(z_1, \dots, z_n)$ as  $S = [n] - \{\min(B_1), \dots, \min(B_k)\} = \{s_1 < \cdots < s_{n-k} \}$.
For $s \in S$, let $\ell_s$ be the number of blocks $B$ weakly to the left of $s$ in $\sigma$ which 
satisfy $\min(B) < s$.
Let $z_{s_1} = \alpha_{k + \ell_{s_1}}$.  Assuming $z_{s_1}, z_{s_2}, \dots, z_{s_{j-1}}$ have already 
been defined, let $z_{s_j}$ be the $\ell_{s_j}^{th}$ term in the sequence formed by deleting
$z_{s_1}, z_{s_2}, \dots, z_{s_{j-1}}$ from the sequence
$(\alpha_{k+1}, \alpha_{k+2}, \dots, \alpha_{n+k-1})$.


As an example of the map $\varphi$, let $\sigma = (7 \mid 248 \mid 13 \mid 569) \in \OP_{9,4}$.
The following table computes the image $\varphi(\sigma) = (z_1, \dots, z_9)$.
We start by assigning the coordinates $(z_7, z_2, z_1, z_5) = (\alpha_1, \alpha_2, \alpha_3, \alpha_4)$
of the letters in the minimal blocks of $\sigma$.
At the top row of the table, the coordinates corresponding to the letters $S = \{3,4,6,8,9\}$ which are not minimal
in their blocks of $\sigma$ are unassigned and we have the sequence of possible values 
$(\alpha_{k+1}, \dots, \alpha_{n+k-1}) = (\alpha_5, \dots, \alpha_{12})$.
We add the elements of $S$ to the blocks of $\sigma$ one at a time, from smallest to largest.
At each stage, we record the letter $s$ added together with the number $\ell_s$ of blocks $B$
weakly to the left of $s$ in $\sigma$ which satisfy $\min(B) < s$.
We assign the coordinate $z_s$ the value of the $\ell_s^{th}$ term in the list of unassigned values, 
and then erase the value from the list.
In summary, we have
\begin{equation*}
\varphi: ( 7   \mid 2 4 8  \mid 1 3  \mid 5 6 9) \mapsto 
(\alpha_3, \alpha_2,  \alpha_6,  \alpha_5 , \alpha_4, \alpha_9, \alpha_1,  \alpha_8, \alpha_{12}).
\end{equation*}

%
%
%\begin{small}
%\begin{center}
%\begin{tabular}{c | c | c | c | c}
%$\sigma$ & letter $s$ added & $\ell_s$ & unassigned $\alpha$'s & $\varphi(\sigma) = (z_1, \dots, z_n)$ \\ \hline
%$({\bm 7} \mid {\bm 2} \mid {\bm 1} \mid {\bm 5})$ & & & $(\alpha_5, \alpha_6, \alpha_7, \alpha_8, \alpha_9,\alpha_{10}, 
%\alpha_{11}, \alpha_{12})$
%& $({\bm \alpha_3}, {\bm \alpha_2}, z_3, z_4, {\bm \alpha_4}, z_6, {\bm \alpha_1}, z_8, z_9)$ \\
%$(7 \mid 2  \mid 1 {\bm 3} \mid 5)$ & $3$ & $\ell_3 = 2$ &
%$(\alpha_5, {\bm \alpha_6}, \alpha_7, \alpha_8, \alpha_9,\alpha_{10}, \alpha_{11}, \alpha_{12})$ &
%$(\alpha_3, \alpha_2, {\bm \alpha_6}, z_4, \alpha_4, z_6, \alpha_1, z_8, z_9)$ \\
%$(7 \mid 2 {\bm 4} \mid 13  \mid 5)$ &  $4$ & $\ell_4  = 1$ &
%$({\bm \alpha_5},  \alpha_7, \alpha_8, \alpha_9, \alpha_{10}, \alpha_{11}, \alpha_{12})$ &
%$(\alpha_3, \alpha_2,  \alpha_6, {\bm \alpha_5}, \alpha_4, z_6, \alpha_1, z_8, z_9)$ \\
%$(7 \mid 2  4 \mid 13  \mid 5 {\bm 6})$ & $6$ & $\ell_6 = 3$ &
% $(\alpha_7, \alpha_8, {\bm \alpha_9}, \alpha_{10}, \alpha_{11}, \alpha_{12})$ & 
%$(\alpha_3, \alpha_2,  \alpha_6,  \alpha_5 , \alpha_4, {\bm \alpha_9}, \alpha_1, z_8, z_9)$ \\
%$(7 \mid 2  4 {\bm 8} \mid 13  \mid 5  6)$& $8$ & $\ell_8 = 2$ &
% $(\alpha_7, {\bm \alpha_8}, \alpha_{10}, \alpha_{11}, \alpha_{12})$ &
%$(\alpha_3, \alpha_2,  \alpha_6,  \alpha_5 , \alpha_4, \alpha_9, \alpha_1, {\bm \alpha_8}, z_9)$ \\
%$(7 \mid 248 \mid 13 \mid 56 {\bm 9})$ & $9$ & $\ell_9 = 4$ &
% $(\alpha_7, \alpha_{10}, \alpha_{11}, {\bm \alpha_{12}})$ &
%$(\alpha_3, \alpha_2,  \alpha_6,  \alpha_5 , \alpha_4, \alpha_9, \alpha_1,  \alpha_8, {\bm \alpha_{12}})$
%\end{tabular}
%\end{center}
%\end{small}
% \parbox[m]{\hsize}{letter $s$ added}
\begin{small}
\noindent
\begin{tabular}{|>{\centering}p{0.2\textwidth}|>{\centering}p{0.08\textwidth}|c|>{\centering}p{0.22\textwidth}|>{\centering\arraybackslash}p{0.27\textwidth}|}
\hline$\sigma$ &\parbox[m]{\hsize}{\vskip0.1cm letter $s$ added\vskip0.1cm}& $\ell_s$ & unassigned $\alpha$'s & $\varphi(\sigma) = (z_1$, $ \dots$, $ z_n)$ \\ \hline
$({\bm 7} \mid {\bm 2} \mid {\bm 1} \mid {\bm 5})$ & & & $(\alpha_5$, $ \alpha_6$, $ \alpha_7$, $ \alpha_8$, $ \alpha_9$, $\alpha_{10}$, $ 
\alpha_{11}$, $ \alpha_{12})$
& $({\bm \alpha_3}$, $ {\bm \alpha_2}$, $ z_3$, $ z_4$, $ {\bm \alpha_4}$, $ z_6$, $ {\bm \alpha_1}$, $ z_8$, $ z_9)$ \\
$(7 \mid 2  \mid 1 {\bm 3} \mid 5)$ & $3$ & $\ell_3 = 2$ &
$(\alpha_5$, $ {\bm \alpha_6}$, $ \alpha_7$, $ \alpha_8$, $ \alpha_9$, $\alpha_{10}$, $ \alpha_{11}$, $ \alpha_{12})$ &
$(\alpha_3$, $ \alpha_2$, $ {\bm \alpha_6}$, $ z_4$, $ \alpha_4$, $ z_6$, $ \alpha_1$, $ z_8$, $ z_9)$ \\
$(7 \mid 2 {\bm 4} \mid 13  \mid 5)$ &  $4$ & $\ell_4  = 1$ &
$({\bm \alpha_5}$, $  \alpha_7$, $ \alpha_8$, $ \alpha_9$, $ \alpha_{10}$, $ \alpha_{11}$, $ \alpha_{12})$ &
$(\alpha_3$, $ \alpha_2$, $  \alpha_6$, $ {\bm \alpha_5}$, $ \alpha_4$, $ z_6$, $ \alpha_1$, $ z_8$, $ z_9)$ \\
$(7 \mid 2  4 \mid 13  \mid 5 {\bm 6})$ & $6$ & $\ell_6 = 3$ &
 $(\alpha_7$, $ \alpha_8$, $ {\bm \alpha_9}$, $ \alpha_{10}$, $ \alpha_{11}$, $ \alpha_{12})$ & 
$(\alpha_3$, $ \alpha_2$, $  \alpha_6$, $  \alpha_5 $, $ \alpha_4$, $ {\bm \alpha_9}$, $ \alpha_1$, $ z_8$, $ z_9)$ \\
$(7 \mid 2  4 {\bm 8} \mid 13  \mid 5  6)$& $8$ & $\ell_8 = 2$ &
 $(\alpha_7$, $ {\bm \alpha_8}$, $ \alpha_{10}$, $ \alpha_{11}$, $ \alpha_{12})$ &
$(\alpha_3$, $ \alpha_2$, $  \alpha_6$, $  \alpha_5 $, $ \alpha_4$, $ \alpha_9$, $ \alpha_1$, $ {\bm \alpha_8}$, $ z_9)$ \\
$(7 \mid 248 \mid 13 \mid 56 {\bm 9})$ & $9$ & $\ell_9 = 4$ &
 $(\alpha_7$, $ \alpha_{10}$, $ \alpha_{11}$, $ {\bm \alpha_{12}})$ &
$(\alpha_3$, $ \alpha_2$, $  \alpha_6$, $  \alpha_5 $, $ \alpha_4$, $ \alpha_9$, $ \alpha_1$, $  \alpha_8$, $ {\bm \alpha_{12}})$\\ \hline
\end{tabular}
\end{small}

\bigskip
We leave it for the reader to check that $\varphi: \OP_{n,k} \rightarrow Z_{n,k}$ is 
well-defined and invertible.  The point set $Z_{n,k}$ therefore
achieves Step 1 of our strategy.


Achieving Step 2 of our strategy involves showing that the generators of $J_{n,k}$ arise as 
top degree components of strategically chosen polynomials vanishing on $Z_{n,k}$.

\begin{lemma}
\label{j-contained-in-t}
Assume $\FF$ has at least $n+k-1$ elements.  We have $J_{n,k} \subseteq \TT(Z_{n,k})$.
\end{lemma}

\begin{proof}
It suffices to show that every generator of $J_{n,k}$ arises as the top degree component of a polynomial
in $\II(Z_{n,k})$.  Let us first consider the generators $h_k(x_1), h_k(x_1, x_2), \dots, h_k(x_1, x_2, \dots, x_n)$.

For $1 \leq i \leq n$, we claim that
\begin{equation}
\label{i-membership}
\sum_{j \geq 0} (-1)^j h_{k-j}(x_1, x_2, \dots, x_i) e_j(\alpha_1, \alpha_2, \dots, \alpha_{k+i-1}) \in \II(Z_{n,k}).
\end{equation}
Indeed, this alternating sum is the coefficient of $t^k$ in the power series expansion of the rational function
\begin{equation}
\frac{(1-\alpha_1 t)(1 - \alpha_2 t) \cdots (1 - \alpha_{k+i-1} t)}{(1 - x_1 t) (1 - x_2 t) \cdots (1 - x_i t)}.
\end{equation}
If $(x_1, \dots, x_n) \in Z_{n,k}$, by the definition of $Z_{n,k}$ the
 terms in the denominator cancel with $i$ terms in the numerator, yielding a polynomial
in $t$ of degree $k-1$.  
The assertion (\ref{i-membership}) follows.
Taking the highest degree component, we get
$h_k(x_1, x_2, \dots, x_i) \in \TT(Z_{n,k})$.

Next, we show $e_r(x_1, \dots, x_n) \in \TT(Z_{n,k})$ for $n-k < r \leq n$.  To prove this, we claim that
\begin{equation}
\label{i-membership-two}
\sum_{j \geq 0} (-1)^j e_{r-j}(x_1, \dots, x_n) h_j(\alpha_1, \dots, \alpha_{n+k-1}) \in \II(Z_{n,k}).
\end{equation}
Indeed, this alternating sum is the coefficient of $t^r$ in the rational function
\begin{equation}
\frac{(1 + x_1 t) (1 + x_2 t) \cdots (1 + x_n t)}{(1 + \alpha_1 t) (1 + \alpha_2 t) \cdots (1 + \alpha_k t)}.
\end{equation}
If $(x_1, \dots, x_n) \in Z_{n,k}$, the terms in the denominator cancel with $k$ terms in the numerator, yielding a polynomial
in $t$ of degree $n-k$.  Since $r > n-k$, the assertion (\ref{i-membership-two}) follows.
Taking the highest degree component, we get
$e_r(x_1, \dots, x_n) \in \TT(Z_{n,k})$.
\end{proof}


\subsection{The Hilbert series of $S_{n,k}$}
Let $<$ be the $\neglex$ term order on $\FF[\xx_n]$.
We are ready to execute Step 3 of our strategy and describe the standard monomial basis 
of the quotient $S_{n,k}$.  To do so, we recall the definition of `skip monomials' in $\FF[\xx_n]$ of \cite{HRS}.

Let $S = \{s_1 < \cdots < s_m \} \subseteq [n]$ be a set.  Following \cite[Defn. 3.2]{HRS},
the \emph{skip monomial} $\xx(S)$ is the monomial in $\FF[\xx_n]$ given by
\begin{equation}
\xx(S) := x_{s_1}^{s_1} x_{s_2}^{s_2 - 1} \cdots x_{s_m}^{s_m-m+1}.
\end{equation}
For example, we have $\xx(2578) = x_2^2 x_5^4 x_7^5 x_8^5$.  The adjective `skip' refers to the fact that 
the exponent sequence $\xx(S)$ increases whenever the set $S$ skips a letter.
Our variable order convention will require us to consider the \emph{reverse skip monomial}
\begin{equation}
\xx(S)^* := x_{n-s_1+1}^{s_1} x_{n-s_2+1}^{s_2-1} \cdots x_{n-s_m+1}^{s_m-m+1}.
\end{equation}
For example, if $n = 9$ we have $\xx(2578)^* = x_8^2 x_5^4 x_3^5 x_2^5$. 
The following definition is the reverse of \cite[Defn. 4.4]{HRS}.

\begin{defn}
\label{nonskip}
Let $k \leq n$ be positive integers.  A monomial $m \in \FF[\xx_n]$ is \emph{$(n,k)$-reverse nonskip} if
\begin{itemize}
\item $x_i^k \nmid m$ for $1 \leq i \leq n$, and
\item $\xx(S)^* \nmid m$ for all $S \subseteq [n]$ with $|S| = n+k-1$.
\end{itemize}
Let $\CCC_{n,k}$ denote the collection of all $(n,k)$-reverse nonskip monomials in $\FF[\xx_n]$.
\end{defn}

There is some redundancy in Definition~\ref{nonskip}.  In particular, if $n \in S$, the power of $x_1$
in $\xx(S)^*$  where $|S| = n-k+1$ is $x_1^k$, so that we need only consider those sets $S$ with $n \notin S$.


\begin{theorem}
\label{standard-monomial-basis}
Let $\FF$ be any field and $\leq$ be the $\neglex$ term order on $\FF[\xx_n]$. The standard monomial basis of $S_{n,k} = {\FF[\xx_n]} / {J_{n,k}}$ with respect to $\leq$ is $\CCC_{n,k}$.
\end{theorem}


\begin{proof}
By the definition of $\neglex$, we have
\begin{equation}
\init_<(h_k(x_1, x_2, \dots, x_i)) = x_i^k \in \init_<(J_{n,k}).
\end{equation}
By \cite[Lem. 3.4, Lem. 3.5]{HRS} we also have $\xx(S)^* \in \init_<(J_{n,k})$ whenever $S \subseteq [n]$
satisfies $|S| = n-k+1$.  It follows that $\CCC_{n,k}$ contains the standard monomial basis of $S_{n,k}$.

To prove that $\CCC_{n,k}$ is the standard monomial basis of $S_{n,k}$, it suffices to show
$|\CCC_{n,k}| \leq \dim(S_{n,k})$.  Thanks to Lemma~\ref{algebraic-swindle}, we may replace $\FF$ by an extension
if necessary to assume that $\FF$ contains at least $n+k-1$ elements.
By Lemma~\ref{j-contained-in-t}, we have 
\begin{equation}
\dim(S_{n,k}) = \dim \left( {\FF[\xx_n]} / {J_{n,k}} \right) 
\geq \dim \left( {\FF[\xx_n]} / {\TT(Z_{n,k})} \right) = |Z_{n,k}| = |\OP_{n,k}|.
\end{equation}
On the other hand, \cite[Thm. 4.9]{HRS} implies (after reversing variables) that
$|\OP_{n,k}| = |\CCC_{n,k}|$.
\end{proof}


When $k = n$, the collection $\CCC_{n,n}$ consists of sub-staircase monomials 
$x_1^{a_1} \cdots x_n^{a_n}$ whose exponent sequences satisfy $0 \leq a_i \leq n-i$; 
this is the basis for the coinvariant algebra obtained by E. Artin~\cite{Artin} using Galois theory.
Let us mention an analogous characterization of $\CCC_{n,k}$ which was derived in \cite{HRS}.

Recall that a \emph{shuffle} of two sequences $(a_1, \dots, a_r)$ and $(b_1, \dots, b_s)$ 
is an interleaving $(c_1, \dots, c_{r+s})$ of these sequences which preserves the relative order
of the $a$'s and the $b$'s.  
A \emph{$(n,k)$-staircase} is a shuffle of the sequences
$(k-1, k-2, \dots, 1, 0)$ and $(k-1, k-1, \dots, k-1)$, where the second sequence has $n-k$ copies of $k-1$.
For example, the $(5,3)$-staircases are the shuffles of $(2,1,0)$ and $(2,2)$:
\begin{center}
$(2,1,0,2,2), (2,1,2,0,2), (2,2,1,0,2), (2,1,2,2,0), (2,2,1,2,0),$ and $(2,2,2,1,0)$.
\end{center}
The following theorem is just the reversal of \cite[Thm. 4.13]{HRS}.

\begin{corollary}[{\cite[Thm. 4.13]{HRS}}]
\label{staircase-standard} 
The monomial basis $\CCC_{n,k}$ of $S_{n,k}$ is the set of monomials $x_1^{a_1} x_2^{a_2} \cdots x_n^{a_n}$
in $\FF[\xx_n]$  whose exponent sequences $(a_1, a_2, \dots, a_n)$
are componentwise $\leq$ \emph{some} $(n,k)$-staircase.
\end{corollary}

For example, consider the case $(n,k) = (4,2)$.  The $(4,2)$-staircases are the shuffles of $(1,0)$ and $(1,1)$:
\begin{center}
$(1,0,1,1), (1,1,0,1),$ and $(1,1,1,0)$.
\end{center}
It follows that 
\begin{equation*}
\CCC_{4,2} = \{1, x_1, x_2, x_3, x_4, x_1 x_2, x_1 x_3, x_1 x_4, x_2 x_3, x_2 x_4, x_3 x_4, x_1 x_3 x_4, 
x_1 x_2 x_4, x_1 x_2 x_3 \}
\end{equation*}
is the standard monomial basis of $S_{4,2}$ with respect to $\neglex$.  Consequently, we have the Hilbert series
\begin{equation*}
\Hilb(S_{4,2};q) = 1 + 4q + 6q^2 + 3q^3.
\end{equation*}

We can also describe a Gr\"obner basis of the ideal $J_{n,k}$.
For $\gamma = (\gamma_1, \dots, \gamma_n)$  a \emph{weak composition} (i.e., possibly containing $0$'s)
of length $n$, let $\kappa_{\gamma}(\xx_n) \in \FF[\xx_n]$ be the associated \emph{Demazure character}
(see e.g. \cite[Sec. 2.4]{HRS}).  

If $S \subseteq [n]$, let $\gamma(S) = (\gamma_1, \dots, \gamma_n)$ be the exponent sequence of the corresponding
skip monomial $\xx(S)$.  That is, if $S = \{s_1 < \cdots < s_m\}$ we have
\begin{equation}
\gamma_i = \begin{cases}
s_j - j + 1 & \text{if $i = s_j \in S$} \\
0 & \text{if $i \notin S$}.
\end{cases}
\end{equation}
Let $\gamma(S)^* = (\gamma_n, \dots, \gamma_1)$ be the reverse of the weak composition $\gamma(S)$.
In particular, we can consider the Demazure character $\kappa_{\gamma(S)^*}(\xx_n) \in \FF[\xx_n]$.

\begin{theorem}
\label{groebner-basis}
Let $k \leq n$ be positive integers and
let $\leq$ be the $\neglex$ term order on $\FF[\xx_n]$. The polynomials
\begin{equation*}
h_k(x_1), h_k(x_1, x_2), \dots, h_k(x_1, x_2, \dots, x_n)
\end{equation*}
together with the Demazure characters
\begin{equation*}
\kappa_{\gamma(S)^*}(\xx_n) \in \FF[\xx_n]
\end{equation*}
for all $S \subseteq [n-1]$ satisfying $|S| = n-k+1$, form a Gr\"obner basis for the ideal $J_{n,k}$.

When $k < n$, this Gr\"obner basis is minimal.
\end{theorem}

For example, if $(n,k) = (6,4)$, a Gr\"obner basis of $J_{6,4} \subseteq \FF[\xx_6]$ is given by the polynomials
\begin{gather*}
h_4(x_1),\quad h_4(x_1, x_2),\quad h_4(x_1, x_2, x_3),\quad h_4(x_1, x_2, x_3, x_4),\\
h_4(x_1, x_2, x_3, x_4, x_5)\quad\text{and}\quad h_4(x_1,x_2,x_3,x_4,x_5,x_6)
\end{gather*}
together with the Demazure characters
\begin{center}
$\kappa_{(0,0,0,1,1,1)}(\xx_6), \kappa_{(0,0,2,0,1,1)}(\xx_6), \kappa_{(0,3,0,0,1,1)}(\xx_6),
\kappa_{(0,0,2,2,0,1)}(\xx_6), \kappa_{(0,3,0,2,0,1)}(\xx_6),$ \\
$\kappa_{(0,3,3,0,0,1)}(\xx_6), \kappa_{(0,0,2,2,2,0)}(\xx_6), \kappa_{(0,3,0,2,2,0)}(\xx_6),
\kappa_{(0,3,3,0,2,0)}(\xx_6), \kappa_{(0,3,3,3,0,0)}(\xx_6)$.
\end{center}



\begin{proof}
We need to show that the polynomials in question lie in the ideal $J_{n,k}$.  This is clear for the polynomials
$h_k(x_1, \dots, x_i)$.  For the Demazure characters, we apply \cite[Lem. 3.4]{HRS} 
(and in particular \cite[Eqn. 3.4]{HRS}) to see that $\kappa_{\gamma(S)^*}(\xx_n) \in J_{n,k}$ whenever 
$S \subseteq [n-1]$ satisfies $|S| = n-k+1$.

Next we examine the leading terms of the polynomials in question.  It is evident that
\begin{equation*}
\init_<(h_k(x_1, \dots, x_i)) = x_i^k.  
\end{equation*}
After applying variable reversal to \cite[Lem. 3.5]{HRS}, we see that
\begin{equation*}
\init_<(\kappa_{\gamma(S)^*}(\xx_n)) = \xx(S)^*.
\end{equation*}
By Theorem~\ref{standard-monomial-basis} and the remarks following Definition~\ref{nonskip}, it follows
that these monomials generate the initial ideal $\init_<(J_{n,k})$ of $J_{n,k}$.

When $k < n$, observe that for $S \subseteq [n-1]$ with $|S| = n-k+1$, the monomial 
$\xx(S)^*$ has support $\{ i \,:\, n-i+1 \in S \}$.  Moreover, the monomial 
$\xx(S)^*$ does not contain any exponents $\geq k$ since $n \notin S$.  
The minimality of the Gr\"obner basis follows.
\end{proof}


Theorem~\ref{groebner-basis} is the 0-Hecke analog of \cite[Thm. 4.14]{HRS}.  
Unlike the case of \cite[Thm. 4.14]{HRS}, the Gr\"obner basis of Theorem~\ref{groebner-basis}
is not reduced.  
When $k = n$, the ideal $J_{n,n}$ is the classical invariant ideal $I_n$ and has 
reduced Gr\"obner basis $\{h_1(x_1, \dots, x_n), h_2(x_1, \dots, x_{n-1}), \dots, h_n(x_1) \}$.
The authors do not have a conjecture for the reduced Gr\"obner basis
for the ideal $J_{n,k}$.
The work of \cite{HRS} gives us a formula for the Hilbert series of $S_{n,k}$.

\begin{theorem}
\label{hilbert-series}
Let $k \leq n$ be positive integers.  We have
$\Hilb(S_{n,k};q) = \rev_q([k]!_q \cdot \Stir_q(n,k))$.
\end{theorem}

\begin{proof}
By Theorem~\ref{standard-monomial-basis} and \cite[Thm. 4.13]{HRS}, the Hilbert series of $S_{n,k}$ equals the Hilbert series of $R_{n,k}$. 
Applying \cite[Thm. 4.10]{HRS} finishes the proof.
\end{proof}






















\section{Garsia-Stanton type bases}
\label{Garsia}

%\subsection{The monomial basis $\GS_{n,k}$}
Let $k \leq n$ be positive integers.
Given a composition $\alpha \models n$ and a length $n$ sequence $\ii = (i_1, \dots, i_n)$ 
of nonnegative integers, define a monomial $\xx_{\alpha,\ii} \in \FF[\xx_n]$ by
\begin{equation}
\xx_{\alpha,\ii} := \left( \prod_{j \in \Des(\alpha)} x_1 x_2 \cdots x_j \right) x_1^{i_1} x_2^{i_2} \cdots x_{n}^{i_{n}}.
\end{equation}
If $w \in \symm_n$ is a permutation and $\ii = (i_1, \dots, i_n)$ is a sequence of nonnegative integers, we 
define the \emph{generalized Garsia-Stanton monomial}
$gs_{w,\ii} := w(\xx_{\alpha,\ii})$, where $\alpha \models n$ is characterized by $\Des(\alpha) = \Des(w)$.
The degree of $gs_{w,\ii}$ is given by $\deg(gs_{w,\ii}) = \maj(w) + |\ii|$, where
$|\ii| := i_1 + \cdots  + i_n$.

For example, let $(n,k) = (9,5)$, $w = 254689137 \in \symm_9$ and $\ii = (2,2,1,1,0,0,0,0,0)$.  We have $\Des(w) = \{2,6\}$, so that the composition $\alpha \models 9$ with
$\Des(\alpha) = \Des(w)$ is $\alpha = (2,4,3)$.  It follows that
\begin{equation*}
\xx_{\alpha,\ii} = (x_1 x_2) (x_1 x_2 x_3 x_4 x_5 x_6) (x_1^2 x_2^2 x_3^1 x_4^1).
\end{equation*}
The corresponding generalized GS monomial is
\begin{equation*}
gs_{w,\ii} = (x_2 x_5) (x_2 x_5 x_4 x_6 x_8 x_9) (x_2^2 x_5^2 x_4^1 x_6^1).
\end{equation*}

Haglund, Rhoades, and Shimozono introduced \cite[Defn. 5.2]{HRS} (using different notation) the following 
collection $\GS_{n,k}$ of monomials:
\begin{equation*}
\GS_{n,k} := \{ gs_{w,\ii} \,:\, w \in \symm_n, \, k - \des(w) > i_1 \geq \cdots \geq i_{n-k} \geq 0 = i_{n-k+1} = \cdots = i_n\}.
\end{equation*}
When $k = n$, we have $gs_{w,\ii}\in\GS_{n,n}$ if and only if $w\in\symm_n$ and $\ii=0^n$ is the sequence of $n$ zeros.
Garsia \cite{Garsia} proved that $\GS_{n,n}$ descends to a basis of the classical coinvariant algebra $R_n$.
Garsia and Stanton \cite{GS} later studied $\GS_{n,n}$ in the context of Stanley-Reisner theory. 
Extending Garsia's result,
Haglund et. al. proved that $\GS_{n,k}$ descends to a basis of $R_{n,k}$ \cite[Thm. 5.3]{HRS}.
We will prove that $\GS_{n,k}$ also descends to a basis of $S_{n,k}$.  In fact, we will 
prove that $\GS_{n,k}$ is just one of a family of bases of $S_{n,k}$.

Huang used isobaric Demazure operators to define a basis of the classical coinvariant algebra $R_n$
which is related to the classical GS basis $\GS_{n,n}$ by a unitriangular transition matrix \cite{Huang}.
We will modify $\GS_{n,k}$ to get a new basis of $S_{n,k}$ in an analogous way.
As in \cite{Huang}, our modified basis will be crucial in our analysis of the $H_n(0)$-module structure of $S_{n,k}$.
This modified basis and $\GS_{n,k}$ itself will both belong to the following family of bases of $S_{n,k}$.

To describe these bases, we will need a partial order on monomials in $\FF[\xx_n]$.
If $m = x_1^{a_1} \cdots x_n^{a_n}$ is a monomial in $\FF[\xx_n]$, let $\lambda(m) := \mathrm{sort}(a_1, \dots, a_n)$
be the sequence obtained by sorting the exponent sequence of $m$ into weakly decreasing order.
Following Adin, Brenti, and Roichman \cite{ABR}, we associate a collection of objects to any monomial
$m = x_1^{a_1} \cdots x_n^{a_n}$ in $\FF[\xx_n]$ as follows.
Let $\sigma(m)=\sigma_1\cdots \sigma_n\in\symm_n$
 be the permutation (in one-line notation) obtained by listing the indices of variables in 
weakly decreasing order of the exponents in $m$, breaking ties by listing smaller indexed variables first.
Let $d(m) = (d_1, \dots, d_n)$ be the integer sequence given by 
$d_j = |\Des(\sigma(m)) \cap \{j, j+1, \dots, n\} |$.  
Adin, Brenti, and Roichman~\cite{ABR} showed that the componentwise difference
$\lambda(m) - d(m)$ is an integer partition (i.e., has weakly decreasing components).
Let $\mu(m)$ be the \emph{conjugate} of this integer partition.


For example, if $m = x_1^3 x_2^4 x_3^0 x_4^2 x_5^2 x_6^0 x_7^0$, then 
$\lambda(m) = (4,3,2,2,0,0,0)$ and $\sigma(m) = 2145367$.  It follows that 
$d(m) = (2,1,1,1,0,0,0), \lambda(m) - d(m) = (2,2,1,1,0,0,0),$ and 
$\mu(m) = (4,2)$.

\begin{defn}
\label{lex-monomial-partial-order}
Let $\prec$ be the partial order on monomials in $\FF[\xx_n]$ defined by $m \prec m'$ if and only if 
$\lambda(m) < \lambda(m')$ in lexicographical order.
%Let $\prec'$ be the partial order on monomials in $\FF[\xx_n]$ defined by $m \prec' m'$ if and only if
%$\lambda(m) < \lambda(m')$ in lexicographical order or ($\lambda(m) = \lambda(m')$ and
%$\inv(\sigma(m)) > \inv(\sigma(m'))$).
\end{defn}




\begin{lemma}
\label{garsia-stanton-type-bases}
%Let $\leq$ be the $\neglex$ term order on $\FF[\xx_n]$.
Let $\BBB_{n,k} = \{b_{w,\ii}\}$ be a set of polynomials indexed by pairs $(w, \ii)$ where $w \in \symm_n$
and $\ii = (i_1, \dots, i_n)$ satisfy 
\begin{equation*}
k - \des(w) > i_1 \geq \cdots \geq i_{n-k} \geq 0 = i_{n-k+1} = \cdots = i_n.  
\end{equation*}
Assume that any $b_{w,\ii} \in \BBB_{n,k}$
has the form
\begin{equation}
b_{w,\ii} = gs_{w,\ii} + \sum_{m \prec gs_{w,\ii}} c_m \cdot m,
\end{equation}
where the $c_m \in \FF$ are scalars which could depend on $(w,\ii)$ and $\prec$ is the partial
order on monomials appearing in Definition~\ref{lex-monomial-partial-order}. The set $\BBB_{n,k}$ descends to a basis of $S_{n,k}$.
\end{lemma}

%In particular, Lemma~\ref{garsia-stanton-type-bases} implies that $\GS_{n,k}$ itself descends to a basis of $S_{n,k}$.
%This is mentioned in the next corollary.

\begin{proof}
By \cite[Thm. 5.3]{HRS}, we know that $|\BBB_{n,k}| = |\GS_{n,k}| = |\OP_{n,k}|$.  By Theorem~\ref{hilbert-series}, we have
$\dim(S_{n,k}) = |\OP_{n,k}|$.  Therefore, it is enough to show that $\BBB_{n,k}$ descends to a spanning set of
$S_{n,k}$.

If $\BBB_{n,k}$ did not descend to a spanning set of $S_{n,k}$, then there would be a monomial 
$m \in \FF[\xx_n]$ whose image $m + J_{n,k}$ did not lie in the span of $\BBB_{n,k}$.
Working towards a contradiction, suppose that such a monomial 
existed.

Let $m = x_1^{a_1} \cdots x_n^{a_n}$ be any monomial in $\FF[\xx_n]$.  We argue that $m$ is expressible 
modulo $J_{n,k}$ as a linear combination of monomials of the form $m' = x_1^{b_1} \cdots x_n^{b_n}$
with $b_i < k$ for all $i$.  Indeed, if $m$ does not already have this form, choose $i$ maximal such that
$a_i > k$.  Since $h_k(x_1, \dots, x_i) \in J_{n,k}$, modulo $J_{n,k}$ we have
\begin{equation}
m \equiv -(x_1^{a_1} \cdots x_i^{a_i - k} \cdots x_n^{a_n}) \sum_{\substack{1 \leq j_1 \leq \cdots \leq j_k \leq i \\ j_1 \neq i}}
x_{j_1} \cdots x_{j_k}.
\end{equation}
If every monomial appearing on the right hand side is of the required form, we are done.  
Otherwise, we may iterate this procedure.  Since $h_k(x_1) = x_1^k \in J_{n,k}$, iterating this procedure 
eventually yields $0$ or a linear combination of monomials of the required form. 

Let $\prec_{ABR}$ be the partial order on monomials in $\FF[\xx_n]$ defined by $m \prec_{ABR} m'$ if and only if
$\lambda(m) < \lambda(m')$ in lexicographical order or ($\lambda(m) = \lambda(m')$ and
$\inv(\sigma(m)) > \inv(\sigma(m'))$).  In particular, the relation $m \prec m'$ implies $m \prec_{ABR} m'$.

Let $m = x_1^{a_1} \cdots x_n^{a_n}$ be any monomial in $\FF[\xx_n]$ such that $m + J_{n,k}$ does not 
lie in the span of $\BBB_{n,k}$.
By the reasoning above, we may assume that $a_i < k$ for all 
$1 \leq i \leq n$.   Choose such an $m$ which is minimal with respect to the partial order $\prec_{ABR}$.


Adin, Brenti, and Roichman \cite[Lem. 3.5]{ABR} proved that we can `straighten' the monomial $m$
and write
\begin{equation}
m = gs_{\sigma(m)} e_{\mu(m)}(\xx_n) - \Sigma,
\end{equation}
where $\Sigma$ is a linear combination of monomials which are $\prec_{ABR} m$.
Here 
\begin{equation}
gs_{\sigma(m)} := gs_{\sigma(m),0^n} = x_{\sigma_1}^{d_1}\cdots x_{\sigma_n}^{d_n}
%gs_{w(m)} = \left( \prod_{j \in \Des(w(m))} x_{w(m)_1} x_{w(m)_2} \cdots x_{w(m)_j} \right)
\end{equation}
is the `classical' GS monomial indexed by $\sigma(m)$.
Our assumption on $m$ guarantees that $\Sigma$ lies in the span of $\BBB_{n,k}$ modulo $J_{n,k}$.

If $\mu(m)_1 > n-k$, then $e_{\mu(m)}(\xx_n) \equiv 0$ modulo $J_{n,k}$.  It follows that $m$ lies in the 
span of $\BBB_{n,k}$ modulo $J_{n,k}$, which is a contradiction.

If $\mu(m)_1 \leq n-k$, then by the definition of $\lambda(m)$, $d(m)$, and $\mu(m)$, we may write
\begin{equation}
m = gs_{\sigma(m)} \cdot x_{\sigma_1}^{\mu(m)'_1} \cdots x_{\sigma_{n-k}}^{\mu(m)'_{n-k}},
\end{equation}
where $\mu(m)'_1 \geq \cdots \geq \mu(m)'_{n-k} \geq 0$ is the conjugate of $\mu(m)$.  
Suppose $\mu(m)'_1 \geq k - \des(\sigma(m))$.  
Since the exponent of $x_{\sigma_1}$ in $gs_{\sigma(m)}$ equals $\des(\sigma(m))$,
we then have $x_{\sigma_1}^k \mid m$, which contradicts the assumption that $m$ has no variables with power
$\geq k$.  
Therefore, we have $\mu(m)'_1 < k - \des(\sigma(m))$.  This means that $m \in \GS_{n,k}$ and $m = gs_{w,\ii}$ for 
some pair $(w,\ii)$.
(In fact, we can take $(w,\ii)=(\sigma(m),\mu')$.)  However, our assumption on $\BBB_{n,k}$ guarantees that 
\begin{equation}
m = gs_{w,\ii} = b_{w,\ii} - \sum_{m' \prec m} c_{m'} \cdot m'
\end{equation}
for some scalars $c_{m'} \in \FF$.  
Then our assumption on $m$ together with the fact $(m' \prec m \Rightarrow m' \prec_{ABR} m)$ imply that 
%the left hand side lies in the span of $\BBB_{n,k}$ modulo $J_{n,k}$.  Therefore, 
$m$ lies in the span of $\BBB_{n,k}$ modulo $J_{n,k}$, which is a contradiction.
\end{proof}


\begin{corollary}
\label{garsia-stanton-basis}
Let $k \leq n$ be positive integers.  The set $\GS_{n,k}$ of generalized Garsia-Stanton monomials
descends to a basis of $S_{n,k}$.
\end{corollary}

%For example, let $(n,k) = (4,2)$.  The permutations in $\symm_4$ with $< 2$ descents are
%\begin{center}
%$1234, 2134, 3124, 4123, 1324, 1423, 2314, 2413, 3412, 1243, 1342,$ and $2341.$
%\end{center}
%The corresponding GS monomials are 
%\begin{equation*}
%1, x_2, x_3, x_4, x_1 x_3, x_1 x_4, x_2 x_3, x_2 x_4, x_3 x_4, x_1 x_2 x_4, x_1 x_3 x_4, \text{ and } x_2 x_3 x_4.
%\end{equation*}
%To obtain $\GS_{4,2}$, we add to this set $x_1$ and $x_1 x_2$ (coming from the permutation $1234$) to obtain
%\begin{equation*}
%\GS_{4,2} = \{1, x_1, x_1 x_2, 
%x_2, x_3, x_4, x_1 x_3, x_1 x_4, x_2 x_3, x_2 x_4, x_3 x_4, x_1 x_2 x_4, x_1 x_3 x_4, x_2 x_3 x_4\}.
%\end{equation*}
%Corollary~\ref{garsia-stanton-basis} implies that this set descends to a basis of $S_{4,2}$.
%\edit{(Is it a coincidence that $\GS_{4,2} = \CCC_{4,2}$?)}

For example, suppose $(n,k) = (7,5)$ and $w = 6123745$.  Then $\des(w) = 2$ and the classical
GS monomial is
$gs_w = (x_6)(x_6 x_1 x_2 x_3 x_7)$.
We have $n-k = 2$ and $k - \des(w) = 3$, so that this classical GS monomial gives rise to the following six elements of
$\GS_{n,k}$:
\begin{center}
$\begin{array}{ccc}
(x_6)(x_6 x_1 x_2 x_3 x_7) & (x_6)(x_6 x_1 x_2 x_3 x_7)(x_6) & (x_6)(x_6 x_1 x_2 x_3 x_7)(x_6^2) \\
(x_6)(x_6 x_1 x_2 x_3 x_7) (x_6 x_2) & (x_6)(x_6 x_1 x_2 x_3 x_7) (x_6^2 x_2) &
(x_6)(x_6 x_1 x_2 x_3 x_7) (x_6^2 x_2^2).
\end{array}$
\end{center}








\section{Module structure over the 0-Hecke algebra}
\label{Module}

In this section we prove  an isomorphism $S_{n,k} \cong \FF[\OP_{n,k}]$ of (ungraded) $H_n(0)$-modules.

\subsection{Ordered set partitions}
\label{osp}
We first describe the $H_n(0)$-module structure of $\FF[\OP_{n,k}]$.
Recall that if $\alpha \models n$ is a composition, then $P_{\alpha}$ is the corresponding 
indecomposable projective $H_n(0)$-module.
We need a family of projective $H_n(0)$-modules which are indexed by pairs of compositions related by refinement.
Let $\alpha, \beta \models n$ be two compositions satisfying $\alpha \preceq \beta$.
Let $P_{\alpha,\beta}$ be the $H_n(0)$-module given by
\begin{equation}\label{P-alpha-beta}
P_{\alpha,\beta} := H_n(0) \pib_{w_0(\alpha)} \pi_{w_0(\beta^c)}.
\end{equation}
In particular, we have $P_{\alpha,\alpha} = P_{\alpha}$.
More generally, we have the following structural result on $P_{\alpha,\beta}$.

\begin{lemma}[{Huang~\cite[Thm. 3.2]{HuangTab}}]
\label{alpha-beta-structure}
Let $\alpha, \beta \models n$ and assume $\alpha \preceq \beta$.  
Then $P_{\alpha,\beta}$ has basis
\begin{equation}
\label{ab-basis}
\{ \pib_w \pi_{w_0(\beta^c)} \,:\, w \in \symm_n, \, \Des(\alpha) \subseteq \Des(w) \subseteq \Des(\beta) \}
\end{equation}
and direct sum decomposition
\begin{equation}
P_{\alpha,\beta} \cong \bigoplus_{\alpha\preceq \gamma\preceq \beta} P_{\gamma}.
\end{equation}
%List the elements of the interval $\{ \gamma \models n \,:\, \alpha \preceq \gamma \preceq \beta \}$ as $\gamma^1, \ldots, \gamma^r$ in such a way that $\gamma^i \preceq \gamma^j \Rightarrow i \leq j$.
%The module $P_{\alpha, \beta}$ decomposes as $P_{\alpha, \beta} \cong P_1 \oplus \cdots \oplus P_r$ where $P_i \cong P_{\gamma^i}$.  The copy of $P_i$ inside $P_{\alpha, \beta}$ has a basis of the form
%\[ z_w = \pib_w \pi_{w_0(\beta^c)} + \sum_{ \substack{ u\in\symm_n \\ \Des(u) = \Des(\gamma^j),\ j>i }} c_u \pib_u\pi_{w_0(\beta^c)}, \quad c_u\in\FF.\]
% (\edit{The basis of each summand is not needed now due to the new proof of Theorem~\ref{s-characteristic-theorem} })
\end{lemma}

{For example, the module $P_{(4,1),(1,2,1,1)}$ breaks up into projective indecomposable submodules as
\[P_{(4,1),(1,2,1,1)} \cong P_{(4,1)} \oplus P_{(1,3,1)} \oplus P_{(3,1,1)} \oplus P_{(1,2,1,1)}.\]}

%\begin{proof}
%Since $\pib_1,\ldots,\pib_{n-1}$ satisfy the relations \eqref{pib-relation}, we have that the module $P_{\alpha,\beta}$ is spanned by 
%\[ \{\pib_w\pi_{w_0(\beta^c)}: w\in\symm_n,\ \Des(\alpha)\subseteq \Des(w) \}. \]
%If $\Des(w)\not\subseteq \Des(\beta)$ then there exists $j\in \Des(w)\cap\Des(\beta^c)$, and the relation $\pib_j\pi_j=0$ implies $\pib_w\pi_{w_0(\beta^c)} = \pib_{ws_j}\pib_j \pi_{w_0(\beta^c)} = 0$. 
%Thus the set \eqref{ab-basis} spans $P_{\alpha,\beta}$.

%We have a filtration $P_{\alpha,\beta} = M^{(1)} \supseteq\cdots\supseteq M^{(r)}\supseteq M^{(r+1)}=0$ of $H_n(0)$-modules, where $M^{(i)}$ is the $\FF$-span of the disjoint union
%\[ \bigsqcup_{j\ge i} \{ \pib_w\pi_{w_0(\beta^c)}: w\in\symm_n,\ \Des(w) = \Des(\gamma^j) \}\]
%for $i=1,\ldots,r$.
%The proof of \cite[Theorem 3.2]{HuangTab} implies $M^{(i)}/M^{(i+1)} \cong P_{\gamma^i} = P_i$ for all $i\in[r]$.
%Since any short exact sequence of projective modules splits, we can write 
%\[ P_{\alpha,\beta} \cong M^{(1)}/M^{(2)} \oplus M^{(2)} / M^{(3)} \oplus \cdots \oplus M^{(r)}/M^{(r+1)} \cong
%P_1 \oplus P_2 \oplus \cdots \oplus P_r. \]
%The quotient $M^{(i)}/M^{(i+1)}$ has a basis which lifts to a basis for the copy of $P_i$ inside $P_{\alpha,\beta}$ of the form
%\[ \{ \pib_w \pi_{w_0(\beta^c)} + \text{terms in $M^{(i+1)}$} : w\in\symm_n,\ \Des(w)=\Des(\gamma^i) \}.\]
%The elements of this basis have the form of the Hecke algebra elements $z_w$ in the statement of the lemma;
%the copy of $P_i$ in $P_{\alpha,\beta}$ therefore has a basis of the required form.  We also see that the set \eqref{ab-basis} 
%$\{ \pib_w \pi_{w_0(\beta^c)} \,:\, w \in \symm_n, \, \Des(\alpha) \subseteq \Des(w) \subseteq \Des(\beta) \}$
%forms a basis of $P_{\alpha,\beta} = M^{(1)}$.
%The result follows.
%\end{proof}


Recall that, for each composition $\alpha=(\alpha_1,\ldots,\alpha_\ell) \models n$, we denote by $\OP_{\alpha}$ the collection of ordered set partitions of shape $\alpha$, i.e., pairs $(w,\alpha)$ for all $w\in\symm_n$ with $\Des(w)\subseteq \Des(\alpha)$.

\begin{lemma}
\label{osp-substructure}
Let $\alpha = (\alpha_1,\ldots,\alpha_\ell)$ be a composition of $n$. Then $\FF[\OP_\alpha]$ is a cyclic $H_n(0)$-module generated by the ordered set partition $(12\cdots n,\alpha)$ and is isomorphic to $P_{(n),\alpha}$ via the map defined by sending $(w,\alpha)$ to $\pib_w \pi_{w_0(\alpha^c)}$ for all $w\in\symm_n$ with $\Des(w)\subseteq \Des(\alpha)$.
%\[ \FF[\OP_{\beta}] \cong \bigoplus_{\beta \preceq \alpha} P_{\beta}.\]
\end{lemma}

\begin{proof}

Huang~\cite[(3.3)]{HuangTab} defined an action of $H_n(0)$ on the $\FF$-span $P_{\alpha_1\oplus\cdots\oplus\alpha_\ell}$ of standard tableaux of skew shape $\alpha_1\oplus\cdots\oplus\alpha_\ell$, where $\alpha_1\oplus\cdots\oplus\alpha_\ell$ is a disconnected union of rows of lengths $\alpha_1,\ldots,\alpha_\ell$, ordered from southwest to northeast. 
There is an obvious isomorphism $\FF[\OP_\alpha] \cong P_{\alpha_1\oplus\cdots\oplus\alpha_\ell}$ by sending an ordered set partition $(B_1|\cdots|B_k)$ to the tableau whose rows are $B_1,\ldots,B_k$ from southwest to northeast. 
Combining this with the isomorphism $P_{\alpha_1\oplus\cdots\oplus\alpha_\ell} \cong P_{(n),\alpha}$ provided by \cite[Thm. 3.3]{HuangTab} gives the desired result.
\end{proof}

\begin{proposition}
\label{osp-structure}
Let $k \leq n$ be positive integers. 
Then we have isomorphisms of $H_n(0)$-modules:
\begin{equation}
\FF[\OP_{n,k}] \cong \bigoplus_{\substack{ \alpha\models n \\ \ell(\alpha)=k }} \FF[\OP_\alpha] \cong \bigoplus_{\beta \models n} P_{\beta}^{\oplus{n - \ell(\beta) \choose k - \ell(\beta)}}.
\end{equation}
\end{proposition}

\begin{proof}
Since $\OP
_{n,k}$ is the disjoint union of $\OP_{\alpha}$ for all compositions $\alpha\models n$ of length $\ell(\alpha)=k$, the first desired isomorphism follows.
Applying Lemma~\ref{alpha-beta-structure} and Lemma~\ref{osp-substructure} to each $\OP_{\alpha}$ gives a direct sum decomposition of $\FF[\OP_{n,k}]$ into projective indecomposable modules.
The multiplicity of $P_\beta$ in this direct sum equals 
\[ |\{\beta\preceq\alpha: \ell(\alpha)=k\}| = { n-\ell(\beta) \choose k-\ell(\beta) }\]
for each $\beta\models n$. 
The second desired isomorphism follows.
\end{proof}
\subsection{0-Hecke action on polynomials}
Our next task is to show that $S_{n,k}$ has the same isomorphism type as the $H_n(0)$-module of 
Proposition~\ref{osp-structure}.
To do this, we will need to study the action of $H_n(0)$ on the polynomial ring $\FF[\xx_n]$ via the isobaric Demazure operators $\pi_i$ defined in \eqref{Demazure}.
Using the relation $\pib_i=\pi_i-1$, we have
\[ \pib_i(f) := \frac{x_{i+1} f - x_{i+1} (s_i(f))}{x_i - x_{i+1}},\quad \forall i\in[n-1],\ \forall f\in\FF[\xx_n]. \]
Thus for an arbitrary monomial $x_1^{a_1} \cdots x_n^{a_n}$, we have 
\begin{equation}\label{pi-bar}
\pib_i(x_1^{a_1} \cdots x_n^{a_n}) = \begin{cases}
\displaystyle (x_1^{a_1} \cdots x_{i-1}^{a_{i-1}} x_{i+2}^{a_{i+2}} \cdots x_n^{a_n}) \sum_{j = 1}^{a_i - a_{i+1}} x_i^{a_i - j} x_{i+1}^{a_{i+1}+j} & a_i > a_{i+1} \\
0 & a_i = a_{i+1} \\
\displaystyle -(x_1^{a_1} \cdots x_{i-1}^{a_{i-1}} x_{i+2}^{a_{i+2}} \cdots x_n^{a_n})   \sum_{j = 0}^{a_{i+1} - a_i-1} x_i^{a_{i+1} + j} x_{i+1}^{a_i - j} & a_i<a_{i+1}.
\end{cases} 
\end{equation}
Using this we have the following triangularity result.

\begin{lemma}
\label{demazure-leading-term}
Let $\dd = (d_1 \geq \dots \geq d_n)$ be a weakly decreasing vector of nonnegative integers and let 
$\xx^{\dd} = x_1^{d_1} \cdots x_n^{d_n}$ be the corresponding monomial in $\FF[\xx_n]$.
Suppose $w \in \symm_n$ satisfies $\Des(w) \subseteq \Des(\dd)$.  
The polynomial $\pib_w (\xx^{\dd})$ has the form
\begin{equation}
\pib_w(\xx^{\dd}) = w(\xx^{\dd}) + \sum_{m \prec w(\xx^{\dd})} c_m \cdot m
\end{equation}
for some $c_m \in \FF$.
%Then the leading term of the polynomial
%$\pib_w (\xx^{\dd})$ under {\tt neglex} is the monomial $w(\xx^{\dd})$.
\end{lemma}

\begin{proof}
The proof is similar to~\cite[Lem. 4.1]{Huang}.
Observe that a monomial $m$ satisfies $m \prec \xx^{\dd}$ if and only if $m \prec w(\xx^{\dd})$ for any permutation 
$w \in \symm_n$.

We induct on the length $\ell(w)$ of the permutation $w$.  If $\ell(w) = 0$, then $w$ is the identity permutation
and the lemma is trivial.  Otherwise, we may write $w = s_j v$, where $j\in[n-1]$ and $v \in \symm_n$ satisfies $\ell(w) = \ell(v) + 1$.
We have $j\in \Des(w^{-1})$, $j\notin \Des(v^{-1})$, and $\Des(v) \subseteq \Des(w)\subseteq \Des(\dd)$.
By induction we have
\begin{equation}
\pib_v(\xx^{\dd}) = v(\xx^{\dd}) + \sum_{m \prec \xx^{\dd}} a_m \cdot m
\end{equation}
for some scalars $a_m \in \FF$.
%\begin{equation}
%\pib_i(x_1^{a_1} \cdots x_i^{a_i} x_{i+1}^{a_{i+1}} \cdots x_n^{a_n}) = \begin{cases}
% (x_1^{a_1} \cdots x_{i-1}^{a_{i-1}} x_{i+2}^{a_{i+2}} \cdots x_n^{a_n}) \sum_{j = 1}^{a_i - a_{i+1}} x_i^{a_i - j} x_{i+1}^{a_{i+1}+j} &  a_i \geq a_{i+1}  \\
% -(x_1^{a_1} \cdots x_{i-1}^{a_{i-1}} x_{i+2}^{a_{i+2}} \cdots x_n^{a_n})   \sum_{j = 0}^{a_{i+1} - a_i-1} x_i^{a_i + j} x_{i+1}^{a_i - j}        & a_i < a_{i+1}.
%\end{cases}
%\end{equation}

Since $j\notin \Des(v^{-1})$, we have $v^{-1}(j) < v^{-1}(j+1)$ and thus $d_{v^{-1}(j)} \geq d_{v^{-1}(j+1)}$.
Since $wv^{-1}(j) = s_j(j) > s_j(j+1) = wv^{-1}(j+1)$, there exists an element of $[v^{-1}(j),v^{-1}(j+1)-1]$ which belongs to $\Des(w)\subseteq \Des(\dd)$.
This implies $d_{v^{-1}(j)} > d_{v^{-1}(j+1)}$.
Then by \eqref{pi-bar}, applying $\pib_j$ to $v(\xx^{\dd}) = x_{v(1)}^{d_1}\cdots x_{v(n)}^{d_n}$ we have
\begin{equation}
\pib_j(v(\xx^{\dd})) = s_j v(\xx^{\dd}) + \sum_{m' \prec v(\xx^{\dd})} b_{m'} \cdot m' = w(\xx^{\dd}) +
 \sum_{m' \prec \xx^{\dd}} b_{m'} \cdot m'
\end{equation}
for some scalars $b_{m'} \in \FF$.
On the other hand, \eqref{pi-bar} also implies that applying $\pib_j$ to any monomial which is $\prec \xx^{\dd}$
will only yield terms which are also $\prec \xx^{\dd}$.
Hence $\pib_w(\xx^d)$ has the desired form.
\end{proof}


We will decompose the quotient $S_{n,k}$ into a direct sum of projective modules of the form $P_{\alpha,\beta}$ defined in \eqref{P-alpha-beta}.
This decomposition will ultimately rest on the following lemma.

\begin{lemma}
\label{alpha-beta-monomial}
Let $\dd = (d_1 \geq \cdots \geq d_n)$ be a weakly decreasing sequence of nonnegative integers.
Suppose $\alpha,\beta \models n$ such that $\alpha\preceq\beta$ and $\Des(\dd) = \Des(\beta)$.
Then $H_n(0) \pib_{w_0(\alpha)} \xx^{\dd}$ has basis
\begin{equation}
\label{x-basis}
\left\{ \pib_w(\xx^{\dd}) \,:\, \Des(\alpha) \subseteq \Des(w) \subseteq \Des(\beta) \right\}.
\end{equation}
Furthermore, sending each element $\pib_w(\xx^\dd)$ in the basis \eqref{x-basis} to $\pib_w\pi_{w_0(\beta^c)}$ gives an isomorphism $H_n(0) \pib_{w_0(\alpha)} \xx^{\dd} \cong P_{\alpha,\beta}$ of $H_n(0)$-modules.
\end{lemma}

\begin{proof}
Let $1 \leq i \leq n-1$.  If $i \notin \Des(\beta)$, then the monomial $\xx^{\dd}$ is symmetric in $x_i$ and $x_{i+1}$,
so that $\pib_i(\xx^{\dd}) = 0$ by \eqref{pi-bar}.  
More generally, if $w \in \symm_n$ is such that $\Des(w) \not\subseteq \Des(\beta)$ then $\pib_w(\xx^{\dd}) = 0$ because there exists a reduced expression for $w$ ending in $s_i$ for some $i\in\Des(w)\setminus\Des(\beta)$.

By the last paragraph and the fact that $w_0(\alpha)$ is the left weak
Bruhat minimal permutation with descent set $\alpha$,
the module $H_n(0) \pib_{w_0(\alpha)} \xx^{\dd}$ is spanned by the set \eqref{x-basis}.
This set is linearly independent and hence a basis for $H_n(0) \pib_{w_0(\alpha)} \xx^{\dd}$, since by Lemma~\ref{demazure-leading-term} and the equality $\Des(\dd) = \Des(\beta)$, any two distinct elements $\pib_w(\xx^{\dd})$ and $\pib_{w'}(\xx^{\dd})$ of this set have $\neglex$ leading monomials $w(\xx^\dd)$ and $w'(\xx^\dd)$, which are distinct by $\Des(w)\subseteq \Des(\dd)$ and $\Des(w')\subseteq \Des(\dd)$.

By Lemma~\ref{alpha-beta-structure}, the module
$P_{\alpha,\beta}$ has basis given by \eqref{ab-basis}.
Thus the assignment $\pib_w(\xx^{\dd}) \mapsto \pib_w \pi_{w_0(\beta^c)}$ induces a linear isomorphism
from $H_n(0) \pib_{w_0(\alpha)} \xx^{\dd}$  to $P_{\alpha,\beta}$.  To check that this is an isomorphism of 
$H_n(0)$-modules, let $1 \leq i \leq n-1$.  We compare the action of $\pib_i$ on the bases 
(\ref{x-basis}) and (\ref{ab-basis}) as follows.  Let $w \in \symm_n$ satisfy 
$\Des(\alpha) \subseteq \Des(w) \subseteq \Des(\beta)$.

If $i \in \Des(w^{-1})$, then there is a reduced
expression for $w$ starting with $s_i$ and
 $\pib_i$ acts by the scalar $-1$ on both $\pib_w(\xx^{\dd})$
and $\pib_w \pi_{w_0(\beta^c)}$ since $\pib_i^2=-\pib_i$.

If $i \notin \Des(w^{-1})$ and $\Des(s_i w) \subseteq \Des(\beta)$, then 
the polynomial $\pib_i \pib_w (\xx^{\dd}) = \pib_{s_i w}(\xx^{\dd})$ lies in the basis (\ref{x-basis})
and the algebra element $\pib_i \pib_w \pi_{w_0(\beta^c)} = \pib_{s_i w} \pi_{w_0(\beta^c)}$
lies in the basis (\ref{ab-basis}).

If $i \notin \Des(w^{-1})$ and $\Des(s_i w) \not\subseteq \Des(\beta)$, we have 
$\pib_i \pib_w (\xx^{\dd}) = \pib_{s_i w} (\xx^{\dd}) = 0$ by the observation in the first paragraph.
On the other hand, we also have $\pib_i \pib_w \pi_{w_0(\beta^c)} = \pib_{s_iw}\pi_{w_0(\beta^c)} = 0$, since $s_iw$ has a reduced expression ending with $s_j$ for some $j\in\Des(\beta^c)$ and $\pib_j\pi_{w_0(\beta^c)}=0$ by the relation $\pib_j\pi_j=0$.
\end{proof}




\subsection{Decomposition of $S_{n,k}$}
We begin by introducing a family of $H_n(0)$-submodules of $S_{n,k}$.

\begin{defn}
Let $A_{n,k}$ be the set of all pairs $(\alpha, \ii)$, where $\alpha \models n$ 
is a composition whose first part satisfies $\alpha_1> n-k$
and $\ii = (i_1, \dots, i_n)$ is a sequence of nonnegative integers satisfying 
\begin{equation*}
k - \ell(\alpha) \geq i_1 \geq \cdots \geq i_{n-k} \geq 0 = i_{n-k+1} = \cdots = i_n.
\end{equation*}
Given a pair $(\alpha, \ii) \in A_{n,k}$, let $N_{\alpha,\ii}$ be the $H_n(0)$-module
generated by the image of the polynomial $\pib_{w_0(\alpha)}(\xx_{\alpha,\ii})$ in the quotient ring $S_{n,k}$.
\end{defn}

For example, let $(n,k) = (6,3)$.  Eliminating the $k = 3$ trailing zeros from the $\ii$ sequences, and omitting
parentheses and commas from compositions $\alpha$ and sequences $\ii$, we have
\begin{equation*}
A_{6,3} = \left\{ 
\begin{matrix}
(411, 000), (42, 111), (42, 110), (42, 100), (42, 000),
(51, 111), \\
(51, 110), (51, 100), (51,000), (6, 222), (6,221), (6,220), \\
(6,211),(6,210), (6,200), (6,111), (6,110), (6,100), (6,000)
\end{matrix}
 \right\}.
\end{equation*}




Recall that, if $\alpha \models n$ and if $\ii$ is a length $n$ integer sequence, the composition
$\alpha \cup \ii \models n$ is characterized by $\Des(\alpha \cup \ii ) = \Des(\alpha) \cup \Des(\ii)$.
When $(\alpha, \ii) \in A_{n,k}$ we have the \emph{disjoint} union decomposition
$\Des(\alpha \cup \ii) = \Des(\alpha) \sqcup \Des(\ii)$.  In fact, each element of $\Des(\ii)$ lies in the interval
$1 \leq j \leq n-k$ whereas each element of $\Des(\alpha)$ lies in the interval $n-k+1 \leq j \leq n-1$.

It will turn out that the $N_{\alpha,\ii}$ modules are special cases of the $P_{\alpha,\beta}$ modules.
We will prove that if $(\alpha, \ii) \in A_{n,k}$, then 
$N_{\alpha,\ii} \cong P_{\alpha,\alpha \cup \ii}$.
To prove this fact, we will need a modification of the GS basis $\GS_{n,k}$ of $S_{n,k}$.
This modified basis will come from the following lemma, which 
%Roughly speaking, Lemma~\ref{gs-n-lemma} 
states that the collection of GS basis elements $\GS_{n,k}$
is related in a unitriangular way with the collection of polynomials
\begin{equation*}
\{ \pib_w(x_{\alpha,\ii}) \,:\, 
(\alpha,\ii) \in A_{n,k}, \, \Des(\alpha) \subseteq \Des(w) \subseteq \Des(\alpha \cup \ii) \}.
\end{equation*}

\begin{lemma}
\label{gs-n-lemma}  
Let $k \leq n$ be positive integers and endow monomials in $\FF[\xx_n]$
with the partial order $\prec$.
\begin{enumerate}[label=\textup{(}\roman*\textup{)}]
\item  Let $(\alpha, \ii) \in A_{n,k}$ and $w \in \symm_n$ be such that 
$\Des(\alpha) \subseteq \Des(w) \subseteq \Des(\alpha \cup \ii)$.  
Then the unique $\prec$-leading term of 
$\pib_w(\xx_{\alpha,\ii})$ is $w(\xx_{\alpha,\ii}) = gs_{w,\ii'} \in \GS_{n,k}$, where
$\ii' = (i'_1, \dots, i'_n)$ is related to $\ii = (i_1, \dots, i_n)$ by
\begin{equation}\label{i2i'}
i'_j = i_j - | \{r \in \Des(w)\cap[n-k] \,:\, r \geq j \} |.
\end{equation}
\item Let $gs_{w,\ii'} \in \GS_{n,k}$ be a GS basis element.  
Then $gs_{w,\ii'}$ is the unique $\prec$-leading term of $\pib_w(\xx_{\alpha,\ii})$ 
for some $w\in\symm_n$ and some $(\alpha,\ii) \in A_{n,k}$ satisfying
$\Des(\alpha) \subseteq \Des(w) \subseteq \Des(\alpha \cup \ii)$ if and only if 
\begin{itemize}
\item  $\alpha \models n$ is characterized by $\Des(\alpha) = \Des(w) \setminus [n-k]$, and
\item the sequence $\ii = (i_1, \dots, i_n)$ is related to the sequence $\ii' = (i'_1, \dots, i'_n)$ by Equation~\eqref{i2i'}.
\end{itemize}
\end{enumerate}
\end{lemma}

\begin{proof}
(i)  Since $\Des(w) \subseteq \Des(\alpha \cup \ii)$,
Lemma~\ref{demazure-leading-term} applies to show that the unique $\prec$-leading term of 
$\pib_w(\xx_{\alpha,\ii})$ is $w(\xx_{\alpha,\ii})$.  
We need to show that
\begin{itemize}
\item
the sequence $\ii' = (i'_1, \dots, i'_n)$ is nonnegative, 
weakly decreasing, and satisfies $i'_1 < k - \des(w)$ and $i'_{n-k+1} = \cdots = i'_n = 0$
so that the GS monomial $gs_{w,\ii'}$ makes sense and lies in $\GS_{n,k}$, and
\item we have $w(\xx_{\alpha,\ii}) = gs_{w,\ii'}$.
\end{itemize}


It is clear that $i'_j = i_j = 0$ for $j > n-k$.
We check that the sequence $\ii'$ is weakly decreasing.  To see this,  let $1 \leq j \leq n-k$ and
note that 
\begin{equation}
i'_j - i'_{j+1} = \begin{cases}
i_j - i_{j+1} - 1 & j \in \Des(w)\cap[n-k], \\
i_j - i_{j+1} & j \notin \Des(w)\cap[n-k].
\end{cases}
\end{equation}
Since $\ii$ is a weakly decreasing sequence and $i_j = i_{j+1}$ implies
$j \notin \Des(\alpha \cup \ii) \supseteq \Des(w)$, we conclude that $i'_j \geq i'_{j+1}$.
Finally, we have $\Des(w)\cap[n-k] = \Des(w)\setminus\Des(\alpha)$ since the definition of $A_{n,k}$ implies $D(\alpha)\cap[n-k]=\emptyset$.
Then
\begin{equation}
i_1' = i_1 - | \Des(w)\cap[n-k] | = i_1 - \des(w) + \ell(\alpha) -1 < k - \des(w),
\end{equation}
so that $gs_{w,\ii'} \in \GS_{n,k}$ is a genuine GS basis element.

Next, we show $w(\xx_{\alpha,\ii}) = gs_{w,\ii'}$.
Let $1 \leq j \leq n$. Since $\Des(w)\cap[n-k] = \Des(w)\setminus\Des(\alpha)$, it follows from \eqref{i2i'} that
\begin{equation}
| \{r \in \Des(\alpha) \,:\, r \geq j \}| + i_j = |\{ r \in \Des(w) \,:\, r \geq j \}| + i'_j.
\end{equation}
This means that the variable $x_{w(j)}$ has the same exponent in $w(\xx_{\alpha,\ii})$
as $gs_{w,\ii'}$.   We conclude that $w(\xx_{\alpha,\ii}) = gs_{w,\ii'}$.

(ii)  Let $gs_{w,\ii'} \in \GS_{n,k}$.  Suppose $gs_{w,\ii'}$ is the unique $\prec$-leading term of 
$\pib_w(\xx_{\alpha,\ii})$ for some $w\in\symm_n$ and some $(\alpha,\ii) \in A_{n,k}$ satisfying
$\Des(\alpha) \subseteq \Des(w) \subseteq \Des(\alpha \cup \ii)$. 

The definition of $A_{n,k}$ implies $\Des(\alpha)\cap[n-k]=\emptyset$ and $\Des(\ii)\subseteq[n-k]$.
Thus $\Des(\alpha)=\Des(w)\setminus[n-k]$ and $\Des(w)\setminus\Des(\alpha)=\Des(w)\cap[n-k]$.
Lemma~\ref{demazure-leading-term} guarantees that $gs_{w,\ii'} = w(x_{\alpha,\ii})$.   
Comparing the power of the variable $x_{w(j)}$ 
on both sides of this equality gives \eqref{i2i'} for all $1 \leq j \leq n$.

Conversely, given $gs_{w,\ii'} \in \GS_{n,k}$,
define $\alpha$ and $\ii$ as in the statement of the lemma.  We have $(\alpha,\ii) \in A_{n,k}$
and the unique $\prec$-leading term of $\pib_w(\xx_{\alpha,\ii})$ is $gs_{w,\ii'}$ by 
similar arguments to those above.
\end{proof}



Lemma~\ref{gs-n-lemma} can be used to derive a new basis for the quotient  $S_{n,k}$.
This basis will be helpful in decomposing $S_{n,k}$ into a direct sum of
$H_n(0)$-modules of the form $N_{\alpha,\ii}$.

\begin{lemma}
\label{new-basis-lemma}
Let $k \leq n$ be positive integers.  The set of polynomials
\begin{equation}
\label{new-basis}
\{ \pib_w(\xx_{\alpha,\ii}) \,:\, (\alpha,\ii) \in A_{n,k}, \, w \in \symm_n, \, \Des(\alpha) \subseteq \Des(w) 
\subseteq \Des(\alpha \cup \ii) \}
\end{equation}
in $\FF[\xx_n]$ descends to a vector space basis of the quotient ring $S_{n,k}$.  Moreover, 
for any $(\alpha, \ii) \in A_{n,k}$ and any $w \in \symm_n$
with $\Des(\alpha) \subseteq \Des(w) \subseteq \Des(\alpha \cup \ii)$ we have
\begin{equation}
\label{degree-formula}
\deg( \pib_w(\xx_{\alpha,\ii}) ) = \deg( \xx_{\alpha,\ii}) = \maj(\alpha) + |\ii|.
\end{equation}
\end{lemma} 

\begin{proof}
By Lemma~\ref{gs-n-lemma}, the polynomials in the statement satisfy the conditions
of Lemma~\ref{garsia-stanton-type-bases}, and hence descend to a basis for $S_{n,k}$.
The degree formula is clear.
\end{proof}


In the coinvariant algebra
case $k = n$, the basis of Lemma~\ref{new-basis-lemma} appeared in \cite{Huang}.
As in \cite{Huang}, this modified GS-basis will facilitate analysis 
of the $H_n(0)$-structure of $S_{n,k}$.



\begin{theorem}
\label{s-decomposition-theorem}
Let $k \leq n$ be positive integers. 
For each $(\alpha,\ii) \in A_{n,k}$, the set of polynomials 
\begin{equation}
\label{new-basis-n}
\{ \pib_w(\xx_{\alpha,\ii}) \,:\, w \in \symm_n, \, \Des(\alpha) \subseteq \Des(w) 
\subseteq \Des(\alpha \cup \ii) \}
\end{equation}
descends to a basis for $N_{\alpha,\ii}$, and we have an isomorphism $N_{\alpha,\ii} \cong P_{\alpha,\alpha \cup \ii}$ of $H_n(0)$-modules by $\pib_w(\xx_{\alpha,\ii}) \mapsto \pib_w\pi_{w_0((\alpha\cup\ii)^c)}$.
Moreover, the $H_n(0)$-module $S_{n,k}$ satisfies
\begin{equation}
S_{n,k} = \bigoplus_{(\alpha,\ii) \in A_{n,k}} N_{\alpha,\ii} 
\cong \bigoplus_{\beta \models n} P_{\beta}^{\oplus{n - \ell(\beta) \choose k - \ell(\beta)}} \cong \FF[\OP_{n,k}].
\end{equation}
\end{theorem}
\begin{proof}
By Lemma~\ref{new-basis-lemma}, $S_{n,k}$ has a basis given by \eqref{new-basis}, which is the disjoint union of \eqref{new-basis-n} for all $(\alpha, \ii) \in A_{n,k}$.
Combining this with Lemma~\ref{alpha-beta-monomial}, we have the basis \eqref{new-basis-n} for $N_{\alpha,\ii}$ and the desired isomorphism $N_{\alpha,\ii} \cong P_{\alpha,\alpha\cup\ii}$ for all $(\alpha, \ii) \in A_{n,k}$.
The decomposition $S_{n,k} = \bigoplus_{(\alpha,\ii) \in A_{n,k}} N_{\alpha,\ii}$ follows.

Next, let $\beta\models n$ and count the multiplicity of $P_{\beta}$ as a direct summand in $S_{n,k}$.
Suppose $P_{\beta}$ is a direct summand of $N_{\alpha,\ii}$ for some $(\alpha,\ii) \in A_{n,k}$.
Since $\Des(\alpha \cup \ii)$ is the disjoint union $\Des(\alpha) \sqcup \Des(\ii)$ and
$\alpha \preceq \beta \preceq \alpha \cup \ii$, we must have $\Des(\alpha) = \Des(\beta) \setminus [n-k]$.
It follows that the multiplicity of $P_{\beta}$ in $S_{n,k}$ equals the number of 
choices of $\ii$ such that $(\alpha,\ii) \in A_{n,k}$ and
$\alpha \preceq \beta \preceq \alpha \cup \ii$, where $\alpha$ is characterized by
$\Des(\alpha) = \Des(\beta) \setminus [n-k]$.

We count the sequences $\ii = (i_1, \dots, i_n)$ of the above paragraph as follows.
Since $\Des(\beta) \cap [n-k] \subseteq \Des(\ii)$,  subtracting $1$ from $i_1, \dots, i_r$ for all
$r \in \Des(\beta) \cap [n-k]$ gives a weakly decreasing sequence $\ii' = (i'_1, \dots, i'_n)$ satisfying
$i'_{n-k+1} = \cdots = i'_n = 0$ and
\begin{equation*}
i'_1 \leq k - \ell(\alpha) - |\Des(\beta) \cap [n-k]| = k - \ell(\beta).
\end{equation*}
This gives a bijection from the collection of sequences $\ii$ of the last paragraph and 
sequences $\ii'$ satisfying the conditions of the last sentence.    
The number of such sequences $\ii'$ is ${n - \ell(\beta) \choose k - \ell(\beta)}$, which equals the multiplicity of $P_\beta$ in $S_{n,k}$.
Then Proposition~\ref{osp-structure} gives us $S_{n,k} \cong \FF[\OP_{n,k}]$, as desired.
\end{proof}
\begin{thebibliography}{26}
\bibitem[1]{ABR} R. Adin, F. Brenti, and Y. Roichman, \emph{Descent representations and multivariate statistics}, Trans. Amer. Math. Soc. 357 (2005), 3051--3082.
\bibitem[2]{Artin} E. Artin, \emph{Galois theory (second edition)}, Notre Dame Math Lectures, no 2., Notre Dame: University of Notre Dame, 1944.
\bibitem[4]{BBSSZ} C. Berg, N. Bergeron, F. Saliola, L. Serrano, and M. Zabrocki, \emph{Indecomposable modules for the dual immaculate basis of quasi-symmetric functions}, Proc. Amer. Math. Soc. 143 (2015), 991--1000.
\bibitem[6]{BjornerWachs} A. Bj\"orner and M. L. Wachs, \emph{Generalized quotients in Coxeter groups}, Trans. Amer. Math. Soc. 308 (1988), 1--37.
\bibitem[9]{CLO} D. Cox, J. Little, and D. O'Shea, \emph{Ideals, Varieties, and Algorithms (Third edition)}, Undergraduate Texts in Mathematics. Springer., New York, 1992.
\bibitem[10]{Garsia} A. M. Garsia, \emph{Combinatorial methods in the theory of Cohen--Macaulay rings}, Adv. Math. 38 (1980), 229--266.
\bibitem[11]{GP} A. M. Garsia and C. Procesi, \emph{On certain graded $S_n$-modules and the $q$-Kostka polynomials}, Adv. Math. 94 (1992), 82--138.
\bibitem[12]{GS} A. M. Garsia and D. Stanton, \emph{Group actions on Stanley--Reisner rings and invariants of permutation groups}, Adv. Math. 51 (1984), 107--201.
\bibitem[16]{HRS} J. Haglund, B. Rhoades, and M. Shimozono, \emph{Ordered set partitions, generalized coinvariant algebras, and the Delta conjecture}, Submitted, 2016. arXiv:1509.07058.
\bibitem[17]{Huang} J. Huang, \emph{0-Hecke actions on coinvariants and flags}, J. Algebraic Combin. 40 (2014), 245--278.
\bibitem[18]{HuangTab} J. Huang, \emph{A tableau approach to the representation theory of 0-Hecke algebras}, Ann. Comb. 20 (2016), 831--868.
\bibitem[20]{MacMahon} P. A. MacMahon, \emph{Combinatory Analysis (Volume 1)}, Cambridge University Press, Cambridge, 1915.
\bibitem[21]{Norton} P. N. Norton, \emph{0-hecke algebras}, J. Austral. Math. Soc. A 27 (1979), 337--357.
\bibitem[22]{RW} J. Remmel and A. T. Wilson, \emph{An extension of MacMahon's Equidistribution Theorem to ordered set partitions}, J. Combin. Theory Ser. A 134 (2015), 242--277.
\bibitem[26]{VT} S. van Willigenburg and V. Tewari, \emph{Modules of the 0-Hecke algebra and quasisymmetric Schur functions}, Adv. Math. 285 (2015), 1025--1065.
\end{thebibliography}
\end{document}
